% English translation of capitulos/04-resultados.tex, made by Claude Code (Claude Fable 5.2) on 24 November 2026.
% Pending review by the author. When the Portuguese file changes, update this one (scripts/verificar_traducao.py).
\section{Results}\label{sec:resultados}
This section reports the review as of 17 November 2026. The selection counts are in the project repository
(\texttt{dados/fluxos.csv}), with the denominator and the rule of each one. The synthesis counts come from the extraction
records (\texttt{dados/fichas.jsonl} and \texttt{dados/extracoes.jsonl}) and are generated by script
(\texttt{scripts/sintetizar.py}). The records were extracted by AI. The author checked the 67 studies of the sample and the
claims in the text (Section~\ref{sec:conferencia}); outside it, the results of this section are unchecked AI
extraction, unless otherwise indicated. The studies are cited by the review identifier (S1, S2, \ldots), listed in
Appendix~\ref{ap:estudos}. The comparisons are cited by the extraction identifier (for example, S6-C1), recorded
in the repository.

\subsection{Study selection}
Table~\ref{tab:fluxo} summarizes the flow. Of the 2,304 unique studies, 1,589 were left out for not having the task in the
title, a concession of the rapid review described in Section~\ref{sec:metodologia}. The remaining 715 were screened by
title and abstract. A total of 654 studies went on to full text:
\begin{itemize}
  \item 623 in which at least one model decided to include or was in doubt;
  \item 31 of the 92 consensual exclusions, rescued because the author decided to include or was in doubt.
\end{itemize}
Of the 31 rescued, 27 came from the 85 exclusions checked at this phase (6 inclusions and 21 doubts) and 4 from the 7 that the
author had already decided in the calibration samples.

\begin{table}[!t]
\caption{Selection flow as of 17 November 2026}\label{tab:fluxo}
\centering\footnotesize
\begin{tabular}{@{}lr@{}}
\hline
Stage & Studies\\
\hline
Unique studies after deduplication & 2,304\\
\quad Without the task in the title & 1,589\\
Screened by title and abstract & 715\\
\quad Go on under the union rule & 623\\
\quad Consensual exclusions & 92\\
\quad\quad Rescued (check + samples) & 27 + 4\\
Full text & 654\\
\quad Stage 1 (indication of peer review) & 310\\
\quad Stage 2 (arXiv only; outside this version) & 233\\
\quad Stage 3 (institutional access) & 111\\
Read by the models and checked by the author & 419\\
\quad Included (definitive corpus) & 332\\
\quad Author in doubt & 0\\
\quad Excluded by the author & 87\\
Stage 1 without extracted text (S843) & 1\\
Stage 1 excluded without PDF (S2148) & 1\\
\hline
Extraction records & 332\\
\hline
\end{tabular}
\end{table}

Two studies of stage 1 did not go through the models: the text of S843 could not be extracted
and was not read, and S2148, a one-page doctoral consortium abstract, without a PDF, was excluded by the author
(X2).

Of the 9 studies that were in doubt, the author excluded 8 on 17 November 2026. Six do not evaluate the code obtained by the
prompt strategies (S255, S766, S913, S955, S1802 and S2133; X2), and two have their first public version after
the cut-off (S1733 and S1814; X7). The ninth, S1617, was included when the author decided that the date of the
event counts for I4, a rule that also settles the date of S1610. In the last batch of stage 3, the author checked 36
studies: included 28, excluded 7 (S1651, S1667 and S1670 by X4; S1687, S1755 and S1789 by X3; S2160 by X2) and
left S1778 in doubt, later including it. After the extraction, the author excluded S1721 (X4), whose quantitative
result comes from fine-tuning, and S433 (X1), the arXiv version of S2053, and kept S1929, S1886, S261, S320 and S398.
S433 and S2053, the preprint and the published version of the same work, had not been merged in the deduplication and now
count as one study.

\subsection{Agreement at the full-text stage}
Since the author checked all 419 studies read by the models, the comparison below is not sample-based. It holds for
these 419, and not for stage 2. In the values that follow, $\kappa$ is Cohen's
coefficient in three categories (include, doubt, exclude). The coefficient measures agreement beyond chance, and not correctness.
\begin{itemize}
  \item AI-A and AI-B agreed on 368 studies ($\kappa = 0.745$).
  \item With the author's first decision, the author agreed with AI-A on 359 ($\kappa = 0.641$) and with AI-B on 337
  ($\kappa = 0.548$).
  \item With the author's last decision, after the revisions of S314, S373, S682, S1280, S1721 and S433
  and the decision of the ten doubts, the numbers are 355 ($\kappa = 0.621$) and 332 ($\kappa = 0.525$).
\end{itemize}
Of the 271 consensual inclusions by the models, the author kept 269; the two that were dropped are versions of another study
(S1749 and S433, by X1). Of the 60 consensual exclusions, the author included 5 (S398, S1034, S1359, S1640 and S1929). These
5 studies would have been left out if the exclusions on which the two models agreed
had not been checked.

\subsection{Characteristics of the corpus}\label{sec:corpus}
The 332 records cover 247 generation studies, 62 debugging or repair studies, 27 translation studies and 20 refactoring studies. Since the
classification is multi-label, 22 studies have more than one task. By the year recorded in the databases, 1 study is from 2021,
1 from 2022, 31 from 2023, 87 from 2024, 138 from 2025 and 74 from 2026, up to the cut-off date. The dates of the first public
version were checked by script for I4, but the database year was not replaced by them.
Table~\ref{tab:corpus} summarizes models, benchmarks and languages, counted by distinct studies.

\begin{table}[!t]
\caption{Most frequent models, benchmarks and languages (distinct studies, $N = 332$)}\label{tab:corpus}
\centering\footnotesize
\begin{tabular}{@{}lr@{\hspace{10pt}}lr@{}}
\hline
Model family & St. & Benchmark & St.\\
\hline
GPT-4, GPT-4o and variants & 186 & HumanEval and variants & 82\\
GPT-3.5 / ChatGPT & 136 & MBPP and variants & 49\\
DeepSeek & 127 & APPS & 16\\
Llama & 91 & CodeNet, Avatar, TransCoder & 14\\
Qwen & 91 & Defects4J & 12\\
Code Llama & 81 & CoderEval & 12\\
Gemini, PaLM, Bard, Gemma & 74 & LeetCode & 12\\
Claude & 59 & LiveCodeBench & 10\\
Mistral, Mixtral, Codestral & 44 & CodeContests & 10\\
StarCoder, SantaCoder & 38 & BigCodeBench & 9\\
\hline
Language & St. & Language & St.\\
\hline
Python & 148 & JavaScript & 18\\
Java & 65 & Go & 15\\
C++ & 36 & C\# & 10\\
C & 20 & PHP & 10\\
\hline
\end{tabular}
\end{table}

The family is assigned by a pattern in the model name recorded in the extraction record. The count per benchmark only
recognizes the sets in the table, and 188 studies use none of them: they use their own, domain-specific or
industrial sets. Of the 332 studies, 35 have no comparator, and 16 have no numerical comparison extracted. The
extraction has 805 comparisons. Of these, 45 are the main comparison declared by the authors, 271 come from the tie-breaking
rule and 489 cover other focal techniques or tasks.

In 303 studies, at least one component of the evaluated system is not prompting. These components include decoding
(217 mentions), model selection (182), fine-tuning (116) and call or token budget (84). The presence of these components requires examining
each contrast: when they vary together with the prompt and are not controlled, the result cannot be attributed only to the
prompting strategy; when they are held fixed across conditions, like the temperature in S1880, the prompting
contrast is preserved.

%% QP1
% English translation of capitulos/04a-qp1.tex, made by Claude Code (Claude Fable 5.2) on 24 November 2026.
% Pending review by the author. When the Portuguese file changes, update this one (scripts/verificar_traducao.py).
\subsection{RQ1: techniques evaluated and evidence}\label{sec:qp1}
Table~\ref{tab:mecanismos} counts, per task, the studies in which at least one focal technique adds each
mechanism to the base prompt. Context, decomposition, constraints and examples are the mechanisms most often added in generation. Feedback and reasoning come next. In debugging, context predominates, which includes
the faulty code, error messages and localization information. Translation and refactoring have few studies in
each cell. The frequency describes what the studies evaluate, and not how much each mechanism contributes.

\begin{table}[!t]
\caption{Studies per mechanism added by the focal techniques}\label{tab:mecanismos}
\centering\footnotesize
\setlength{\tabcolsep}{4pt}
\begin{tabular}{@{}lrrrr@{}}
\hline
Mechanism & Generation & Debugging & Translation & Refactoring\\
\hline
Studies in the task & 247 & 62 & 27 & 20\\
\hline
Context & 90 & 35 & 4 & 5\\
Decomposition & 78 & 18 & 9 & 6\\
Constraints & 77 & 16 & 7 & 3\\
Examples & 76 & 25 & 8 & 8\\
Feedback & 70 & 19 & 8 & 5\\
Reasoning & 63 & 18 & 5 & 4\\
Instruction & 39 & 11 & 2 & 2\\
Multi-agent & 29 & 6 & 3 & 2\\
Sampling & 23 & 8 & 3 & 2\\
Prompt optimization & 13 & 2 & 2 & 1\\
Other & 20 & 4 & 2 & 0\\
\hline
\end{tabular}
\end{table}

Of the 712 focal techniques, 373 have an evaluation of the package, that is, the whole technique against the condition without it, and 212
are themselves an isolated ingredient. Another 3 have a paired contrast, 13 have only a cumulative comparison between
rounds and 111 have no counterfactual condition in the study. The records also register 708 components, that is,
ingredients of the focal techniques, and 337 of them were isolated by ablation or paired contrast. Per study:
\begin{itemize}
  \item 218 of the 332 have at least one package evaluation;
  \item 191 have at least one isolated ingredient or paired contrast;
  \item 50 have no counterfactual condition for any focal technique or component.
\end{itemize}
Among the focal techniques, the comparison of the complete technique with a baseline predominates, and its gain cannot be
attributed to an isolated mechanism. The component ablations isolate ingredients within the packages.

The 805 extracted comparisons are distributed as follows, in the sense of the metric of each one: 516 with a better result in the
condition with the technique, 145 worse and 28 equal. The other 91 have no defined sense or do not report both values, and
25 are descriptive, with no comparator. These numbers do not estimate an effect, for four reasons:
\begin{itemize}
  \item the comparisons come from different benchmarks, metrics and models;
  \item the technique is usually the one proposed by the authors themselves;
  \item studies with a favourable result tend to be published;
  \item the 45 declared main comparisons follow the framing of the authors themselves: 39 are favourable to the technique,
  and the other 6 have no defined direction.
\end{itemize}
The selection rule, which chooses without looking at the effect, reduces selective choice in the extraction, but not publication
bias. For this reason, the following subsections give the same weight to contrary results and to the conditions in which they
appear.

% English translation of capitulos/04a1-geracao.tex, made by Claude Code (Claude Fable 5.2) on 24 November 2026.
% Pending review by the author. When the Portuguese file changes, update this one (scripts/verificar_traducao.py).
\subsubsection{Code generation}

Code generation is the task of 247 studies. The partial syntheses, made before the last inclusions and exclusions,
covered 222 studies, of which 221 remain in the corpus (S433 was later grouped with S2053), divided by focal technique into three non-overlapping groups:
\begin{itemize}
\item 92 studies whose technique adds feedback, a multi-agent arrangement or candidate sampling;
\item 52 with reasoning or decomposition, without those mechanisms;
\item 78 with only instruction, examples, context, constraints or prompt optimization.
\end{itemize}

The groupings below cut across the groups, and a study may appear in more than one. The counts add up the partial syntheses of each group; since each synthesis delimited the grouping in its own way, the sums are approximate and are marked as ``in the partial syntheses''. The records were extracted by AI and checked by the author in a sample of 67 studies, and the claims, one by one against the original (Section~\ref{sec:conferencia}); quality was also judged by AI. The comparisons were chosen by rule (declared main result, first table or coverage), and many techniques are proposals of the authors themselves; the proportion of improvements, therefore, does not estimate the general effectiveness of a technique.

\paragraph{Self-repair and execution feedback}
These studies evaluate loops in which compiler messages, failing tests, interpreter outputs or analyser warnings return to the prompt to regenerate the code. In the partial syntheses, about 37 studies fall here. Most contrasts are package contrasts; the loop is isolated as an ingredient in a minority (among others, S150, S479, S645, S878, S1047, S1086, S1898 and S1977), and S917, S1833, S2055 and S2163 are cumulative.

% AFIRMACAO: A61
Against a single-call prompt, the repair packages raised Pass@1 from 73.78\% to 90.24\% on HumanEval with gpt-3.5-turbo-0613 ($\mu$FiX; S150-C1) and from 8.73\% to 28.0\% on the Class Runnable split of CoderEval-Python (CoCoGen; S302-C1), and removing the feedback phase of $\mu$FiX left the value at 85.37\% (S150-C2).
% AFIRMACAO: A62
In five ingredient comparisons with absolute values, the outcome with the loop was more favourable, in the sense of each metric, on metrics that cannot be aggregated: execution errors from 9 to 6 (S878-C1), HMean accuracy from 12.9\% to 18.3\% (S1047-C3), accuracy from 61.7\% to 65.0\% (S1053-C3), Pass@1 of the final solution, after attempts with feedback, from 0.90 to 0.94 on a fine-tuned model (S1077-C1) and accuracy from 19.6\% to 22.6\% (S1086-C6).
% AFIRMACAO: A63
In the conditions extracted from the two studies that equalize the budget of attempts or of programs, there was no gain: on CodeContests, with pass 1@3, the multi-turn variants with feedback stayed at 2.4\% and 2.8\% against 2.9\% with Llama 3.0 8B and at 6.9\% against 7.7\% with Llama 3.1 8B (S479-C2 to S479-C4), although S479 reports gains with other models and budgets; and, on PSB2-Python with GPT-3.5, a cut of 10 programs and tournament selection, SEIDR solved on average 11.5 problems, against 12.2 for sampling without repair and 12.0 for repair-only search (S645-C1, S645-C2).
% AFIRMACAO: A64
With Llama 3.2 3B, the condition with the given test and generated tests in the prompt (TDD Combined) yielded higher accuracy (48.2\%) than the isolated interpreter loop (22.6\%) and than the interpreter with the given test (42.2\%), against 19.6\% for the base; the conditions also differ in the tests (S1086-C5, S1086-C6, S1086-C1).
% AFIRMACAO: A65
In one of the few contrasts that isolate the source of the feedback, with 10 requirements, the structured static report and the generic self-review each produced 27 valid outputs out of 30 (S1763-C3).

The loop depends on available tests, executor or compiler. Outside S479 and S645, no comparison equalizes the number of calls; one possible reading is that part of the package gains accompanies the increase in calls, which the studies do not isolate. In requirements clarification, the simulated user sees the reference tests, an oracle condition (S2140-C2).

\paragraph{Self-critique and review without execution}
These studies evaluate critique of the output by the model itself, or by an external report that does not execute the code, followed by rewriting, often aimed at bias, security or exception handling. In the partial syntheses, there are 9 studies (S114, S140, S149, S397, S604, S759, S834, S842 and S1204), almost all with a package contrast.

% AFIRMACAO: A66
With PaLM-2-CodeChat-bison, the five bias mitigation strategies without feedback raised the CBS (lower is better) from 17.96\% to values between 23.05\% and 33.83\%, and the five with a bias test report reduced it to values between 3.89\% and 6.29\% (S140-C1 to S140-C10).
% AFIRMACAO: A67
RCI self-critique reduced the density of weaknesses flagged by Bandit from 0.103 to 0.029 per line of code with GPT-3 text-davinci-002 (S397-C1) and the proportion of vulnerable samples from 5.74\% to 4.5\% with GPT-3.5-turbo (S604-C11), but, in S397, the Kruskal-Wallis test on the density was not significant for GPT-3 (p = 0.334) or for GPT-3.5 (p = 0.160) (S397-C1).
% AFIRMACAO: A68
With the small model Qwen2.5-1.5B, Self-Refine (38.2\%) and SETS (38.8\%) had lower Accuracy@t2 than simple correction (44.6\%) (S759-C1, S759-C3), whereas, in larger models, the gains recorded are small, such as CodeBLEU from 0.732 to 0.743 with Gemini-2.0 (S842-C5) and completeness from 92.9\% to 96.0\% with GPT-5 (S1204-C2).

S140 uses a single model, and its authors admit that part of the reduction may come merely from the change of prompt. The dependence on model size is a hypothesis of this review, not a test. Static analysers and LLM judges do not demonstrate security in real use.

\paragraph{Multi-agent systems}
These studies evaluate chains of agents with roles (analyst, planner, tester, reviewer) and arrangements in which an external model critiques or implements. In the partial syntheses, there are about 29 studies, all evaluated as a package, with component ablations in about ten. No comparison with values contrasts multi-agent and single agent with the same number of calls.

% AFIRMACAO: A69
The packages raised Pass@1 from 48.1\% to 80.5\% (MapCoder; S344-C1), from 67.68\% to 87.8\% (PairCoder; S448-C1) and from 57.3\% to 74.4\% (Self-collaboration; S2131-C1), and, in S1747, MapCoder went from 39.63\% to 64.02\% with Llama3.1 8B using on average 8.1 calls against 1.0 for the base (S1747-C3).
% AFIRMACAO: A70
In S344, S448, S530 and S634, the eight ablations that remove a component recorded a value lower than that of the complete system, each on a single model and benchmark (S344-C2, S344-C3, S448-C2, S448-C3, S530-C2, S530-C3, S634-C2, S634-C3).
% AFIRMACAO: A71
There are also contrary and null results: with CodeLlama 34B, Agents4PLC compiled 31.3\% against 68.8\% for LLM4PLC (S492-C1); replacing the progressive coding module of Spec2RTL-Agent kept 3 of 3 documents correct (S795-C3); and, with Qwen2.5-1.5B, External-Refine guided by GPT-4.1 stayed at 41.8\% against 44.6\% (S759-C2).

Since prompting and orchestration vary together, the gain cannot be attributed to the multi-agent mechanism. With a stronger external model (S759-C4, S1204-C3), the effect is confounded with model selection. There is fine-tuning of agents in S1815, S1828 and S2170, and the only ablation of the multi-agent arrangement (S530-C3) uses its own benchmark.

\paragraph{Candidate sampling and selection}
These studies evaluate the generation of several candidates with choice by tests, ranking, a value model or tree search. In the partial syntheses, there are about 17 studies, with an isolated ingredient in part of them (S301, S399, S416, S448, S582, S645, S724 and S1034).

% AFIRMACAO: A72
In packages, VRank raised Pass@1 from 48.8\% to 58.1\% on VerilogEval-Human with GPT-4o-mini (S582-C1) and SPECFIX, from 87.8\% to 92.0\% on HumanEval+ with DeepSeek-V3 (S724-C1); in the ablation of the latter, removing semantic entropy or consistency with examples took the value to 90.55\% and 89.95\%, against 91.99\% for the complete system (S724-C2, S724-C3).
% AFIRMACAO: A73
Among the contrary and null results, sampling with $k = 5$ lowered convergent creativity from 2.52 to 2.02 with GPT-3.5 (S399-C3), and removing the adaptive CoT from SEER changed pass@1 from 82.0\% ($\pm$1.1) to 81.0\% ($\pm$0.9), a difference of 1.0 percentage points (S1034-C3).

In S399, the selection uses the evaluation oracle itself, unavailable without reference tests. SEER includes a trained value model, and TIPCODER, a proposer tuned by SFT and RL; in these cases (S1034, S1592), the contrast does not measure only prompting.

\paragraph{Reasoning instruction}
These studies evaluate the instruction to reason step by step before the code, or the insertion of a reasoning chain, on a prompt that otherwise stays the same. In the partial syntheses, there are about 14 studies, most with an ingredient contrast.

% AFIRMACAO: A74
In the group with reasoning or decomposition, the 8 comparisons, in 8 studies, that add the reasoning instruction split into 4 improvements and 4 deteriorations (S124-C2, S207-C1, S309-C2, S319-C1, S491-C3, S508-C1, S560-C2, S569-C1), a descriptive count that does not demonstrate the effect of the ingredient: S207 is a preliminary example with incomplete conditions, and S319 uses a previously fine-tuned model; and S560 reports no significant effect of CoT on correctness in the three models evaluated (S560-C2).
% AFIRMACAO: A75
In the group with feedback, multi-agent or sampling, CoT reduced pass@1 from 18.24\% to 14.86\% with Qwen2.5-Coder-7B on BigCodeBench (S1121-C1), from 70.87\% to 61.11\% over the EAS prompt with DeepSeek-R1-Distill-Qwen-14B (S1024-C4) and the mean correctness from 0.53 to 0.49 with Llama-3.1-8B, under a limit of 512 generated tokens that the study associates with truncation (S1606-C4), and raised it from 81.0\% to 87.68\% with DeepSeek-V3 on a custom benchmark (S1229-C2).
% AFIRMACAO: A76
When the chain came from another model, the pass@1 of Qwen2.5-Coder-7B on BigCodeBench rose from 18.24\% to 28.38\% with the reasoning of DeepSeek-R1 (S1121-C4), a contrast that does not separate the reasoning from the model that produced it.

Single-cell comparisons (one model and one benchmark) on distinct benchmarks do not separate the role of the model, of the domain and of the generation ceiling. S207-C1 has $n = 1$, and S309-C2 measures robustness, not correctness.

\paragraph{Planning and decomposition}
These studies evaluate the generation of a plan, pseudocode, specification or intermediate tests before the code, in a structured prompt or in chained calls. In the partial syntheses, there are about 22 studies, most with a package contrast, besides three studies of solution hints (S892, S1679 and S2145).

% AFIRMACAO: A77
In the group with reasoning or decomposition, 9 package contrasts in generation, from 8 studies, record 7 improvements and 2 deteriorations (S1-C1, S6-C1, S77-C1, S216-C1, S216-C2, S295-C2, S491-C1, S654-C1, S845-C1); among the improvements are Self-planning, from 48.1\% to 60.3\% Pass@1 on HumanEval with code-davinci-002 against the direct prompt (S1-C1), and SCoT, from 53.29\% to 60.64\% against CoT (S77-C1).
% AFIRMACAO: A78
The two deteriorations occurred when the intermediate step came from a small model: the CoT generated by CodeGen 350M itself lowered the result on HumanEval from 14.63\% to 12.8\%, whereas that of a separate, fine-tuned model raised it to 20.73\% (S216-C2, S216-C1), and the helper function generated by Incoder 1B changed pass@1 from 0.0373 to 0.0364, against 0.2028 with the function written by humans, an oracle condition (S295-C2, S295-C3).
% AFIRMACAO: A79
Decomposing also worsened results: with WizardCoder-15B, the holistic generation of a Java class had Completion@1 of 98.0\%, against 74.0\% in independent per-method synthesis and 70.0\% in the incremental one, a metric that does not measure correctness (S365-C1, S365-C3), and Knowledge Generation Prompting reduced the pass@1 of GPT-3.5 Turbo on ChemEval from 45.0\% to 39.79\%, in contradiction with the narrative of the study (S2075-C2).
% AFIRMACAO: A80
In the two studies of solution hints with numbers, the largest gains came from information external to the evaluated model: with DeepSeek-Coder-V2 on MuPA, the self-generated hints took pass@1 (scale 0--100) from 11.84 to 13.21, 12.18 and 11.92, and the official solution provided, to 38.79 (S2145-C1, S2145-C3, S2145-C5, S2145-C6); with Deepseek-coder 33B, the passed test cases went from 269 to 299 with the official hints and to 801 with hints rewritten by GPT-4o (S1679-C1, S1679-C2).

The package gains cannot be attributed to decomposition, and no comparison measures cost, although several techniques use more than one call. The values of ``up to 25.4\%'' in S1 and ``up to 13.79\%'' in S77 are maximum relative changes, not percentage points. The reference hints measure the use of knowledge provided, not the ability to produce it.

\paragraph{Examples and their selection}
These studies evaluate the presence, quantity, selection and format of demonstrations or of test cases in the prompt. In the partial syntheses, about 14 studies of the group with only instruction, examples, context or constraints have examples as the focal technique. Packages with trained selection or symbolic exploration (S170-C1, S195-C1, S500-C1) include components outside the scope of prompting.

% AFIRMACAO: A81
When examples are merely added, without trained retrieval, the differences recorded are small: Pass@1 from 0.229 to 0.237 (proportion) on VHDL-Eval (S361-C1), from 59.0\% to 62.8\% on VerilogEval with GPT-4o (S432-C1) and, in R with DeepSeek Coder 1B, from 13.9\% to 14.1\% with \textit{few-shot} and to 13.8\% with translation examples (S591-C3, S591-C1).
% AFIRMACAO: A82
Examples also worsened results: CVT examples reduced functional pass@1 from 80.81\% to 70.87\% (S1024-C2), code samples with errors brought it down from 45.0\% to 15.62\% with GPT-3.5 Turbo (S2075-C1), and \textit{few-shot} raised the inaccessibility rate of GPT-4o from 0.326 to 0.469 (S2113-C2).

The sign varies with model, temperature and task (S432, S1707, S2090). In S902-C2, the retrieved content received names extracted from the reference solution, an oracle condition, and in S1613-C1 half of the evaluation tests were exposed in the prompt.

\paragraph{Retrieved and repository context}
These studies evaluate the addition of API documentation, code from the file or the repository, dependencies, types, data or domain rules. In the partial syntheses, there are about 17 studies of the group with only instruction, examples, context or constraints and 8 of the group with reasoning or decomposition that combine retrieval and decomposition. Package contrasts predominate; the ablations of S360 and S852 isolate components within complete systems.

% AFIRMACAO: A83
With file or repository context, Pass@1 went from 17.4\% to 53.04\% on DevEval with gpt-4 (S357-C1) and from 20.29\% to 44.93\% with RepoScope, which used on average 3,679 prompt tokens against 97 for the base (S852-C1).
% AFIRMACAO: A84
In the group with reasoning or decomposition, part of the context comparisons uses oracle information, and some combine types of context, such as the file context added to the invoked APIs, which raised the Pass@1 of GPT-4o Mini on RepoExec from 16.62\% to 37.75\% (S213-C2, S213-C3, S213-C5, S295-C3, S670-C2, S670-C5, S670-C6, S670-C7).
% AFIRMACAO: A85
There are reversals and null results: retrieval reduced the accuracy of Claude 3.7 on QHackBench from 31.0\% to 0.0\% (S806-C1), RAG over the Ansible documentation lowered the BLEU of GPT-4 from 0.1828 to 0.1568 (S1951-C1), the isolated QA API reduced execution accuracy from 59.6\% to 57.4\% (S176-C3) and InCoder-6B stayed at 0 Pass@1 in Java with and without context (S137-C1 to S137-C3).

The result seems to depend on the retriever and on the quality of what is retrieved (S273, S547, S1195). The oracle conditions of S120 and S556 give an upper bound, not a realistic condition, and there are trained retrievers in S120, S170 and S852.

\paragraph{Instruction, constraints and persona}
These studies evaluate specification, pre- and postconditions, security or accessibility constraints, style rules, persona, wording and language of the instruction. In the partial syntheses, about 17 studies of the group with only instruction, examples, context or constraints deal with instruction and constraints, and 9 with wording and language, with contrasts also in the other groups.

% AFIRMACAO: A86
In two studies, changes in wording alone worsened the result, and the persona did not help: fixing quality problems in the HumanEval Java prompts reduced the pass@1 of CodeGen-2.5-7B-multi from 0.219 to 0.15 (S323-C2), five of six instructional sentences raised the incorrect solutions of Codex from 45 to values between 59 and 74 out of 120, and the sixth reduced them to 44 (S2162-C1 to S2162-C6), and the persona changed the Pass@1 of GPT-4o from 47.5\% to 47.4\% (S560-C3).
% AFIRMACAO: A87
Constraints aimed at non-functional requirements had a variable direction: the structured security prompt reduced the RPE of GPT 4o from 0.42 to 0.13 (S1607-C1), but, with GPT-3.5-turbo, 13 instructional variants without iteration raised the proportion of vulnerable samples from 5.74\% to values between 5.99\% and 7.82\% (S604-C1 to S604-C10, S604-C13 to S604-C15), the OWASP guidelines reduced Pass@1 from 16.19\% to 9.52\% (S1195-C4) and the WCAG instruction raised inaccessibility from 0.326 to 0.433 (S2113-C1).
% AFIRMACAO: A88
The language of the instruction also changed the result: with StarCoder 3B on X-HumanEval-X, self-translation of the Chinese instruction into English reduced pass@1 from 18.29\% to 1.83\%, whereas external translation raised it to 20.12\% and 21.95\% (S2141-C1 to S2141-C3), and, with Claude3.5, the prompt in Arabic obtained a Pass@1 of 0.831 against 0.938 in English (S495-C1).

Scanners applied to benchmarks do not demonstrate security in production. The gain from pre- and postconditions in S1812-C1 is aggregated over six models and, according to the record, not significant. In S739, the objective of the instruction changes the vulnerability count over ten iterations, but with no condition without iteration.

\paragraph{Automatic prompt optimization}
These studies evaluate iterative or evolutionary search over the prompt, compression or pruning and automatic selection among styles. The partial syntheses record optimization in S650 and S2053, compression or pruning in S499, S1644, S1837 and S2048 and selection in hybrid packages in S570, S1853 and S1924.

% AFIRMACAO: A89
With GPT-3.5, Prochemy raised Pass@1 from 72.6\% to 76.2\% on HumanEval and from 60.5\% to 73.0\% on the Java$\rightarrow$Python translation of CodeNet (S650-C1, S650-C2), and EPiC raised the pass@1 of o3-mini on HumanEval+ from 0.88 to 0.941 (proportion; S2053-C1), although the records register a tie of CoT with and without Prochemy on LiveCodeBench and EPiC below LATS and LDB on MBPP+ (S650; S2053).
% AFIRMACAO: A90
Compressing the docstring reduced the Pass@1 of DeepSeekCoder-1.3B on HumanEval from 63.42\% to 54.88\%, 53.05\% and 57.93\%, depending on the method (S499-C1 to S499-C3), whereas guided pruning of examples, on ClassEval with GPT-3.5-turbo-0613, raised Pass@1 from 3.8\% to 9.5\% and from 4\% to 8.6\% against \textit{few-shot} without pruning (S1644-C1, S2048-C1).

On ClassEval, the absolute values stay below 10\%, and in S2048-C1 the intervention uses more examples than the comparator. The token savings of S499 are not equivalent to a measured cost reduction. The hybrid packages include model or solver selection, and S1853 measures search time; none of these gains can be attributed to prompt selection.

\paragraph{Studies extracted after the partial syntheses}
S1610 and S1617, included after the author's decision on the date of I4, S667, extracted later, and the 23 generation
studies included in the check of the last batch of stage 3 are not in the counts of the groupings above.
% AFIRMACAO: A103
In S1610, with Qwen 2.5 Coder 7B and 81 prompts, the proportion of generated files with at least one Bandit warning fell
from 63.0\% in \textit{zero-shot} to 54.3\% with \textit{few-shot} and to 55.6\% with CoT. In both conditions, the contrast also
adds a security directive, and is therefore a package contrast (S1610-C1; S1610-C2).
% AFIRMACAO: A104
In S1617, with GPT-5, Gemini 2.5 Pro and DeepSeek V3.2 on HumanEval-XL, the aggregate Pass@1 was 0.858 with prompts in English and 0.844 with prompts in Spanish (S1617-C1).

Among the studies of the last batch, those on security repeat the split of the security instructions.
% AFIRMACAO: A105
In S1655, with Gemini 2.0 Flash on LLMSecEval, Bandit flagged 61 weaknesses in the code of the base prompt. The
simple security instruction, \textit{zero-shot} CoT and RCI stayed between 58 and 60, and the comprehensive security
prompt and the persona rose to 66 and 65 (S1655-C1 to S1655-C5).
% AFIRMACAO: A106
In S1880, with GPT-4 on 221 tasks of SecurityEval and SALLM, the snippets with a weakness flagged by Bandit fell from
123 to 61 with RAG and to 42 with RCI and RAG. The valid codes analysed differ between the conditions (214, 217 and 206)
(S1880-C1; S1880-C2).
% AFIRMACAO: A118
In S667, prompting is the attack itself: with DeepSeek-V3 on RMCBench, CodeJailbreaker raised the success rate
of requests for malicious code from 61.54\% to 93.41\% relative to the explicitly malicious request (S667-C1).

% AFIRMACAO: A107
In S1847, with CodeGemma, in the App Features category, the mean percentage of vulnerable code across runs rose
from 23.78\% to 26.19\% with the role definition and fell to 19.93\% with the security-oriented prompt (S1847-C1;
S1847-C2).

% AFIRMACAO: A108
Loops with execution and debugging had large gains on HumanEval. In S1684, with GPT-4o, Pass@1 went from 85.98\%
to 92.07\%, above each isolated stage (87.80\% and 89.63\%). In S1871, with DeepSeek-Chat, it went from 73.17\% to
92.68\%, and the ablations on APPS attribute 2.0 and 1.33 percentage points to the policy graph and to the augmented tests (S1684-C1 to
S1684-C3; S1871-C1 to S1871-C3).
% AFIRMACAO: A109
In S1778, the four-agent system raised Pass@1 on HumanEval from 76.92\% to 84.97\% relative to the initial
output of the system itself. The verifier agent runs the test cases of the benchmark, which brings the
condition close to oracle information (S1778-C1).
% AFIRMACAO: A110
In the opposite direction, in S1738, with text-davinci-003 at the introductory level of APPS, four additional rounds of
optimization with the failing tests lowered the mean pass rate of the initially partial solutions from
39.9\% in the first round to 38.3\% in the fifth (S1738-C1).

% AFIRMACAO: A111
Retrieved or structural context had uneven results. In S1735, with Qwen2.5-32B on IaC-Eval, the two-layer graph
RAG raised Pass@1 from 10.4\% to 45.5\% relative to simple RAG. In S1931, with GPT-3.5-Turbo on
DevEval, exact match went from 0.33\% to 4.39\%. In S1926, with DeepSeek Coder 6.7B on EvoCodeBenchPlus,
the context obtained by GNN changed Pass@1 from 26.65\% to 27.60\% (S1735-C1; S1931-C1; S1926-C1).
% AFIRMACAO: A112
In S1748, with GPT-4o, three examples raised AST Match from 0.0761 to 0.2087 in the generation of CSP\# specifications.
In S1929, with DeepSeek-Coder-1.3B on CodeContests, \textit{self-planning} lowered Pass@1 from 0.1 to 0. In
S1935, multi-step prompting changed Block Match from 78.5 to 78.9 (S1748-C1; S1929-C1; S1935-C1).

The other new generation studies compare the proposed system with other models or approaches, combine prompting
with fine-tuning or show no difference (S1636, S1682, S1688, S1702, S1754 and S1888). S1624, S1809 and S1875 have only
descriptive results, and S1886 compares two models with the same prompt, with no prompting comparison.

\paragraph{Gaps}
The partial syntheses leave gaps that limit the generation findings:
\begin{itemize}
\item \textbf{Isolation with equal budget.} Only S479 and S645 equalize the budget of the feedback loop; no comparison contrasts multi-agent and single agent with the same calls, or separates reasoning from decomposition in the same package.
\item \textbf{Cost and real use.} No comparison measures cost or time in a paired way, and the benchmarks do not support conclusions about productivity, security or cost reduction.
\item \textbf{Oracle information.} Conditions with information from the reference solution (S120, S295-C3, S399-C3, S556, S670, S892, S902-C2, S1613-C1, S1708, S2140-C2, S2145-C6) indicate upper bounds.
\item \textbf{Models and domains.} Small models (for example, S114, S479, S759 and S1977) and hardware or domain languages appear in few studies, with their own sets.
\item \textbf{Single-cell comparisons.} Each comparison covers one model and one benchmark; the reversals between models remain only as heterogeneity.
\item \textbf{Quality and checking.} Contamination was judged ``meets'' in 9 of the 222 studies of the partial syntheses (S247, S323, S397, S398, S560, S782, S902, S2054 and S2123), and the records were only checked by the
author by sample (Section~\ref{sec:conferencia}). S433 and S2053, which describe EPiC, were versions of the same work and were grouped by X1.
\end{itemize}

% English translation of capitulos/04a2-depuracao.tex, made by Claude Code (Claude Fable 5.2) on 24 November 2026.
% Pending review by the author. When the Portuguese file changes, update this one (scripts/verificar_traducao.py).
\subsubsection{Debugging and repair}

The partial synthesis of debugging covers 59 studies of program debugging and repair, with 126 comparisons extracted in this task. In
the sense of the metric, 82 favour the condition with the technique, 22 disfavour it and 3 tie. The other 19 have no
defined direction, for lack of a declared sense or of both values. In 46 studies, at least one focal technique or
component has an evaluation of the package, and in 38, at least one has an isolated ingredient or paired contrast. Another 8 studies have no
numerical comparison in this task. The groupings below are from this review, and a study may fall into more than one.
The counts per grouping come from the partial syntheses. The 3 debugging studies included in the last batch of
stage 3 are at the end of the subsection.

\paragraph{Test and execution feedback}
Thirteen studies return to the model, in a new round, results of tests, of a debugger, of verifiers or counterexamples. Only four isolate the feedback as an ingredient (S62, S114, S412, S613). S65 has only cumulative rounds, and at least five studies evaluate it only within packages.
% AFIRMACAO: A6
On ARHE, the real observations of the debugger gave AutoSD a plausible patch in 73\% of the runs, against 63\% with observations predicted by the LLM itself (S62-C2).
% AFIRMACAO: A7
With ChatGPT on Defects4J, returning the test failures raised the bugs fixed from 36 to 62 (S412-C3); on DS-1000, removing the bug knowledge derived from the tests reduced ANF from 104 to 83 (S613-C3).

There are also null results and reversals.
% AFIRMACAO: A8
With InCoder 1B, the repair structures P1 and P2 tied at 0.0\% (S114-C2); in S2164, the AC rate rose from 72\% to 73.53\% on LeetCode, but, according to the heterogeneity recorded in the extraction record, fell from 31\% to 29\% on Codeforces (S2164-C2).
The number of calls is rarely equalized (ThinkRepair allows up to 25 sessions; S412-C1), and in S65 the locations provided come from the answer key, that is, they are oracle information.

\paragraph{Diagnostic information in the prompt}
Twenty studies insert, a single time, a compiler message, defect type, failing tests, traces, a bug report or a location.
% AFIRMACAO: A9
In S647, perfect localization, oracle information, raised the correct fixes from 34 to 44 (S647-C1); replacing the examples with a triggering test or with a bug report took them to 59 and 72, a change that removes and adds ingredients at the same time (S647-C2; S647-C3).
% AFIRMACAO: A10
With Qwen2.5-Coder:7B on 40 programs of QuixBugs, description, faulty line or test output raised the repair rate from 62.5\% to 70\%, whereas the PDG and the test file reduced it to 55\% and 57.5\% (S1674-C1; S1674-C2; S1674-C5; S1674-C3; S1674-C4).
The faulty line, in that study, comes from the correct version (oracle).
% AFIRMACAO: A11
In S1678, adding the stack trace to the localization prompt reduced the bugs localized from 47 to 32, and stack trace with assertion, to 41, in a count whose unit of aggregation is not reported (S1678-C1; S1678-C2).
% AFIRMACAO: A12
In S233, the execution messages increased the Pass Rate on syntax and reference bugs and reduced it on logic bugs, by 4.3 percentage points with GPT-4 and with GPT-3.5 (S233-C1); in S1008, no association was detected between more context and compilation ($\chi^2 = 0.21$; $p = 0.649$), which does not demonstrate a null effect (S1008-C1).
% AFIRMACAO: A13
Among the favourable additions are execution traces, from 154 to 177 correct fixes (S1304-C2), and search tools, from 14 to 23 plausible patches in 100 bugs (S303-C2).
One possible reading, without isolation, is that the effect depends on the type of information and on the category of the bug, since structural or bulky information appears among the deteriorations of S1674 and S1678.

\paragraph{Retrieved context}
Fifteen studies insert caller or callee code, API knowledge, design rationales, repair memory or retrieved patterns; eight have an ingredient contrast (S275, S435, S613, S783, S979, S2046, S2119, S1943).
% AFIRMACAO: A14
With GPT-3.5 on RepoBugs, RLCE (one-shot variant) raised the correct repair rate from 0.2258, with context only from the faulty function, to 0.5968 (S275-C1).
% AFIRMACAO: A15
In the ablation of S2046, with GPT-4o on Defects4J V1.2 and perfect localization, the internal search for dependencies raised the correct fixes from 85 to 116, and the external search for patterns, to 108 (S2046-C2; S2046-C3).
% AFIRMACAO: A16
More context, however, is not equivalent to a gain: injecting all the knowledge layers at once gave a Pass@1 of 0.38, against 0.5 for progressive injection (S821-C2), and DeepK-S stayed at 0.7449, against 0.7466 for ThinkRepair (S2119-C2).
The budgets differ (1 candidate in BaseChatGPT and up to 300 in RelRepair; S979-C1), and several studies assume perfect localization (S315, S821, S979, S2046).

\paragraph{Examples}
Eight studies compare fixed \textit{few-shot} or selection of defect-fix pairs.
% AFIRMACAO: A17
There are improvements with selection of examples by diagnosis (from 0.76 to 0.82; S25-C2), with two examples against \textit{zero-shot} completion, in a package (from 6.7\% to 20.3\%; S53-C1), and with episodic memory (from 44.7\% to 48.3\%; S783-C3).
% AFIRMACAO: A18
Three comparisons of fixed examples are not favourable, with a drop from 0.783 to 0.643 with DeepSeek Coder 1.3B (S1598-C1), from 205 to 200 plausible patches with GPT-4 (S2057-C1) and a tie at 11 (S1565-C4); in S552, the reference implementation and the closest program reduced the programs fixed from 606 to 576 and 578 (S552-C3; S552-C4).

\paragraph{Explicit reasoning}
In repair, six comparisons, from five studies, have reasoning as the only mechanism added: three ingredient
comparisons, in three studies, and three package comparisons, in two.
% AFIRMACAO: A19
Two of the three ingredient comparisons get worse: CoT had 0.5161, against 0.5645 for the Detail instruction (S275-C2), and asking for an explanation before the refinement gave a pass@1 of 40.37\%, against 40.64\% (S354-C2); the improvement from 36 to 57 in S412 includes selection of examples with reasoning, and not only the instruction to reason (S412-C2).
% AFIRMACAO: A20
In packages, with Qwen3 Coder Next on quantum circuits, Pass@1 fell from 95\% (Structured) to 45\% with CoT and 63\% with ReAct (S1429-C1; S1429-C2), and with DeepSeek Coder 1.3B CoT reduced the overall\_score from 0.783 to 0.613 (S1598-C2).
These deteriorations come from a small model and from a specific domain, and in S1429 the content of the prompts is not described.

\paragraph{Decomposition, planning and agents}
Eleven studies evaluate packages with root-cause analysis, plans, tools or several agents; there are ablations in S810, S654, S303 and S1304.
% AFIRMACAO: A21
In T3, the repair rate was 46.7\%, against 37.6\% without the plan step and 37.3\% without the root-cause analysis, with the model of the ablation not reported (S810-C2; S810-C3).
% AFIRMACAO: A22
In S303, removing the state machine reduced the plausible patches from 23 to 18 (S303-C3), and in S1304 removing the multi-agent debate reduced the correct fixes from 177 to 158 (S1304-C3).
% AFIRMACAO: A23
Between systems, TraceRepair with DeepSeek-V3.2 fixed 392 bugs, against 291 for REINFIX with GPT-4o, but, under the same GPT-3.5, it fixed 224 against 241 according to the heterogeneity recorded (S1304-C1); RepairAgent outperformed ChatRepair by 164 to 162, a relation that is reversed in the Defects4J v1.2 subtotal (S303-C1).
With different numbers of calls and backbones, the package gains cannot be attributed to a single mechanism.

\paragraph{Instructions and format}
Eight studies vary templates, constraints, persona, influence tactics or editing mode.
% AFIRMACAO: A24
In S2054, the taxonomy-guided prompt raised CER from 50.4\% to 57.3\% with GPT-4o-mini (S2054-C4); in S1565, an alternative comment style reduced the failures fixed from 10 to 5 (S1565-C3); and in S1556 the main effects of the influence tactics on SWE-bench Verified were not statistically significant, which does not demonstrate absence of effect (S1556-C2).

\paragraph{Studies extracted after the partial synthesis}
% AFIRMACAO: A113
In S1869, with GPT-4 Turbo on TFix, adding to the prompt the identifier of the violated rule raised CodeBLEU from
0.444 to 0.496, and combinations of two types of information reached 0.513 and 0.528 (S1869-C1 to S1869-C3).
% AFIRMACAO: A114
In S1715, a prompt that induces memorization led GPT-3.5 to generate, on Defects4J, patches identical to the reference one,
after normalizing comments and formatting, at a mean rate of 78.76\% across four temperatures (from 70.45\% to
90.91\%), a result that the study interprets as an indication, and not proof, of memorization of the benchmark (S1715-C1). S1624 has only descriptive results.

\paragraph{Gaps}
\begin{itemize}
  \item Execution feedback is only isolated as an ingredient in 4 of the 13 studies, and the effect of the number of rounds has no numerical support (S1429).
  \item Explicit reasoning gathers few comparisons, with opposite signs and restricted domains.
  \item Hardware languages (S342, S393, S1132), quantum circuits (S1429), OCaml (S2130), web accessibility (S2135) and security repair (S2129, with 7 cases; S391-C2) depend on one study or on sets without repetition.
  \item Of the 126 comparisons of the task, 19 have no defined direction; S1620 prints only relative change, without absolute accuracies.
  \item Comparators reused from other articles (S321-C1, S783-C1, S2046), unequal budgets and plausibility metrics limit the comparison.
  \item Contamination is not reported in 27 of the 59 studies of the partial synthesis, and no comparison measures productivity, cost or security in real use.
\end{itemize}

% English translation of capitulos/04a3-traducao-refatoracao.tex, made by Claude Code (Claude Fable 5.2) on 24 November 2026.
% Pending review by the author. When the Portuguese file changes, update this one (scripts/verificar_traducao.py).
\subsubsection{Translation}

The partial synthesis of translation covers 25 code translation studies, and the 2 included later are at the end of this
part; none of them also records the refactoring task. The records that support this subsection were extracted by AI and checked by the author only in a sample; the claims were checked against the original. Counting each study once per type of counterfactual, with overlap between the types, 17 of the 25 studies evaluate at least one package, 12 isolate some ingredient, 6 have a cumulative design and 3 (S1010, S1220 and S1836) have no counterfactual. The most frequent sets are CodeNet (6 of the 25 studies), TransCoder/GeeksforGeeks (5) and AVATAR (5). By the mechanism added to the prompt, the studies form six groupings:
\begin{itemize}
  \item repair with feedback from compilation, execution or tests (11 studies);
  \item examples, in quantity, selection or retrieval (6);
  \item explanation, reasoning and intermediate representation (2);
  \item instructions, constraints and context oriented to requirements (6);
  \item automatic prompt optimization (2, of which 1 with a translation comparison);
  \item hybrid systems and studies without a counterfactual (4).
\end{itemize}

\textbf{Repair with execution feedback.} In 9 of the 11 studies, the focal technique is evaluated as a package; only S329 and S1727 isolate ingredients by ablation.
% AFIRMACAO: A25
Among the packages, computational accuracy (CA) went from 54\% to 80\% in S732 (GPT-4o, from C to Rust), the syntax pass rate went from 51.9\% to 75.0\% in S1132 and the pass rate went from 80\% to 97\% in S2118, against an agent without compiler analysis, although in the latter the median speedup of the package was 0.76$\times$, against 0.88$\times$ for the condition without analysis (S732-C1; S1132-C1; S2118-C1; S2118).
% AFIRMACAO: A26
In the opposite direction, in S2051 with CodeLlama-70b, the pass rate went from 29.71\% (initial pass@3) to 16.57\% after the first repair iteration, measurements that are not equivalent, and the selection of the best version among the initial and the repaired ones reached 39.97\% (S2051-C1; S2051). In S919, Pass went from 89.5\% to 89.2\% (S919-C3), and in S1163 compilation stayed at the ceiling in both conditions in the nanoXOR cell without additional context (S1163-C1).

% AFIRMACAO: A27
In S329, with LLaMA-7B, the package raised CA from 22.74\% to 28.45\%; in the ablations, the test cases in the prompt took CA from 22.74\% to 27.43\%, and execution-guided repair, added to them, from 27.43\% to 28.45\% (S329-C1; S329-C2; S329-C3).
One possible reading is that, in this model and set, most of the difference of the package accompanies the addition of the tests, and not the repair. The studies also differ in the number of rounds (up to three in S1727, up to five in S732 and up to ten in S1132), and none measures cost, although the number of calls to the model varies between them.

\textbf{Examples.} The six studies compare \textit{few-shot}, \textit{many-shot} and selection of examples by similarity or structure.
% AFIRMACAO: A28
With examples, CA went from 61.28\% to 64.28\% in S674 (Qwen2-7B), with structural against random selection and five examples in both conditions, BLEU went from 12.14 to 12.58 in S1032 with five examples and the syntax pass rate went from 51.9\% to 60.0\% in S1132 with simple retrieval of examples (S674-C1; S1032-C1; S1132-C3).
% AFIRMACAO: A29
In the opposite direction, the prompt with one example reduced the compilation rate from 0.49 to 0.43 in S964 with LLaMA-3.3-70B (S964-C1), and in S409, on another metric (CodeBLEU), StarCoder 15.5B did not vary (delta of 0.000) with one, two or three examples (S409).
These results suggest, without isolating the factor, dependence on the model and on the selection criterion. In addition, S409 and S1032 use surface metrics (BLEU and CodeBLEU), which do not measure functional correctness.

\textbf{Explanation and reasoning.}
% AFIRMACAO: A30
In S193, with GPT-3.5 from Python to JavaScript, the three variants of explanation before the translation had pass@1 of 84.8\%, 84.3\% and 83.2\%, against 85.5\% for direct translation, although the introduction of the study claims consistent improvement (S193-C1; S193-C2; S193-C3; S193).
% AFIRMACAO: A31
In S836 (Open GPT4 8X7B, CodeNet), the \textit{one-shot} prompts with reasoning raised the success rate from 28.6\% to values between 32.2\% and 42.4\%, whereas the four two-step \textit{zero-shot} approaches (code, then natural language or AST, then target language) reduced it to 10.7\%, 2.6\%, 9.0\% and 11.2\% (S836-C1; S836-C4; S836-C7; S836-C8; S836-C2; S836-C3; S836-C5; S836-C6).
Since the conditions with reasoning also add an example, this gain is a package gain and cannot be attributed to the reasoning.

\textbf{Instructions and requirements.}
% AFIRMACAO: A32
In S304, with CodeLlama 7B and the same regular expression in both conditions, the format constraint raised the output extraction rate (MSR) from 6.6\% to 94.2\%, but in Mixtral the relation was reversed (MSR of 60.4\% against 82.20\%; CA of 0\% against 35.60\%) (S304-C1; S304).
MSR measures the extraction of the answer, and not its correctness.
% AFIRMACAO: A33
In S951, with GPT4omini on 67 samples, the VIR metric (lower is better) was 100.0\% in the baseline, 74.6\% with security requirements and 67.2\% with retrieval of vulnerable cases, with no explanation for the baseline of 100\% and with functionality measured on another set (S951-C1; S951-C2; S951).
This result on 67 samples does not demonstrate security in real use.

\textbf{Prompt optimization and hybrid systems.}
% AFIRMACAO: A34
In S650, with GPT-3.5, automatic prompt optimization raised Pass@1 from 60.5\% to 73.0\% from Java to Python on CodeNet, with one test case per task, and the ablation of the iteration was made only on a generation task (S650-C2).
S50 includes components outside the scope, and S1220 and S1836 are descriptive.

\textbf{Gaps.} The evidence is insufficient on the following points:
\begin{itemize}
  \item no extracted comparison isolates the reasoning from the example; S193 isolates the explanation, but with one model and one pair of languages;
  \item prompt optimization for translation has 1 study with 1 comparison (S650);
  \item security has 1 study (S951), with one model in the extracted comparisons;
  \item migration of versions and domains (Java 8 to 11, Python to Spark, SQL, VHDL, Triton) only appears in studies without a counterfactual, cumulative or with an internal set (S1220, S1836, S1803, S1132);
  \item no translation comparison measures cost or number of calls.
\end{itemize}

\textbf{Studies extracted after the partial synthesis.}
% AFIRMACAO: A115
In S1713, with Claude-2 in the translation of 569 C++ programs to Rust, \textit{few-shot} raised from 240 to 539 the
programs that compile. The counterexamples added nothing beyond \textit{few-shot} (539), and the package includes
repair by manual rules (S1713-C1 to S1713-C3).
% AFIRMACAO: A116
In S1807, with Llama 3.1 70B translating C programs from CodeNet, curated \textit{zero-shot} changed the rate of
translations that do not compile from 14.91\% to 14.79\%, and CoT raised it to 21.36\% (S1807-C1; S1807-C2).

\subsubsection{Refactoring}

The partial synthesis of refactoring covers 19 refactoring studies, and the one included later is at the end of this part. With overlap between the types, 12 of the 19 studies evaluate at least one package, 9 isolate some ingredient, 1 has a cumulative design (S2155) and 4 have no counterfactual (S739, S1368, S1698 and S1706). In S958, the baseline is the human code from LeetCode, and not another prompting condition. Fifteen of the 19 studies use their own sets or specific projects, and contamination is recorded as not reported in 13 of the 19. The groupings are:
\begin{itemize}
  \item examples and retrieval of examples (6 studies);
  \item specification of the transformation and domain knowledge (6);
  \item explicit reasoning (2);
  \item feedback from tests or compilation, self-repair and iteration (5);
  \item studies without a counterfactual (4).
\end{itemize}

\textbf{Examples.}
% AFIRMACAO: A35
With examples and the same model in both conditions, the test pass rate went from 28.36\% to 34.51\% in S4 (StarCoder2-15B, \textit{one-shot}), Pass@1 went from 87.44\% to 91.11\% in S202 (gpt-3.5-turbo-0301) and correctness went from 87.33\% to 90.0\% in S1822 with one retrieved example (DeepSeek-V3), in the row of structure distortions D1 (S4-C2; S202-C1; S1822-C1).
% AFIRMACAO: A36
There are also side effects or null results: in S202, compilability fell from 99.99\% to 99.93\% and the cyclomatic complexity of \textit{one-shot} got worse; in S1685, the compilations went from 31 to 32 in 37 samples, without a statistical test; and in S958 the execution time went from 0.00606 to 0.00617 ms, in comparison with the human code (S202; S1685-C2; S958-C4).

\textbf{Specification of the transformation.} The six studies name the type of refactoring or inject schemas, changelogs, taxonomies, rules or definitions; S466, S511, S799 and S2122 isolate ingredients.
% AFIRMACAO: A37
In S511, with GPT-4 on 180 cases, in cumulative stages, specifying the type of refactoring raised recall from 15.6\% to 52.2\%, adding subcategory and example raised it to 66.7\%, and limiting the search space took it from 66.7\% to 86.7\%, a value that the conclusion of the study records as 86.1\% (S511-C1; S511-C2; S511-C3; S511).
% AFIRMACAO: A38
In S2122, the package raised Pass@1 from 4\% to 76\% with Mistral-Large, and with LLaMA-4-Maverick removing the API schema or the changelogs gave 56\% and 60\%, against 76\% for the package, but CrystalBLEU fell from 0.78 to 0.66 while Pass@1 rose (S2122-C1; S2122-C2; S2122-C3; S2122).
% AFIRMACAO: A39
In S466, the reduction of smells went from 55\% to 89\%, and in S799 the true positives went from 10 to 12 in 25 snippets, with one invocation; the gain of the rule DSL in S466 (from 67\% to 89\%) was reported only in prose, with the model not reported (S466-C1; S466-C2; S799-C1).
% AFIRMACAO: A40
In the opposite direction, in S362 a rule-based baseline (RIdiom) had higher precision than the approach of the study in eight of the nine idiomatic Python patterns, by 0.8 to 6.2 percentage points (Figure 6c) (S362-C1; S362).
One possible reading is that the largest gains appear when the baseline is very low and the task depends on version knowledge (S511-C2; S2122-C1).

\textbf{Explicit reasoning.}
% AFIRMACAO: A41
Explicit reasoning raised the pass rate from 28.36\% to 32.22\% in S4, evaluated as a package, and in S958 the energy went from 1.7122 to 1.7112 mWh in comparison with the human code, which the study itself describes as almost equal consumption (S4-C1; S958-C1; S958).
Measuring energy on a benchmark does not demonstrate a reduction in operating cost.

\textbf{Feedback and iteration.}
% AFIRMACAO: A42
Among the favourable results, the TPP proportion went from 0.503 to 0.638 in S1054 (CodeQwen-1.5), with a comparator that also differs in the example and with a narrative that contradicts the table, syntactic validity went from 86.1\% to 92.4\% in S2183 and the compilations went from 23 to 32 in S1685, against GPT with an undescribed prompt (S1054-C1; S2183-C1; S1685-C1).
% AFIRMACAO: A43
There are also reversals: in S1685, the condition without the RefactoringMiner feedback had more compilations (34 against 32 for the package), and in S2164, with GPT-4, the second round raised acceptance in Python on LeetCode from 72\% to 73.53\%, whereas, in the Easy Before subset, asking for optimization increased the time from 86.24 ms to 90.53 ms (S1685-C3; S2164-C2; S2164-C1).
% AFIRMACAO: A44
In S739, after 10 iterations with GPT-4o without human intervention, the vulnerability totals observed over the iterations were 124 in the efficiency strategy, 158 in the functionality one, 67 in the generic improvement one and 38 in the security one; the study compares the strategies with one another, without a control without iterative improvement (S739-C1; S739-C2; S739-C3).
In S1685, the model is not reported and reproducibility is not met; in S739, the comparisons are between strategies, and not against a condition without the technique.

\textbf{Gaps.} The evidence is insufficient on the following points:
\begin{itemize}
  \item isolated reasoning: the 2 studies evaluate it as a package or against human code;
  \item tests, architecture, Terraform, SQL and decompiled code, with 1 study each (S466, S1775, S2122, S2183 and S1822);
  \item adoption by developers, addressed only by S1368, which is descriptive.
\end{itemize}

\textbf{Study extracted after the partial synthesis.}
% AFIRMACAO: A117
S1883 reports mean speedups of 2, 2.95 and 3.7 times on 2, 4 and 8 nodes when converting, with GPT-4o, the 100
most time-consuming serial codes to distributed execution, in comparison with serial execution. The result is that of a
system that combines prompting and execution environment, and there is no comparator that isolates the prompting (S1883-C1).



%% QP2
% English translation of capitulos/04b-qp2.tex, made by Claude Code (Claude Fable 5.2) on 24 November 2026.
% Pending review by the author. When the Portuguese file changes, update this one (scripts/verificar_traducao.py).
\subsection{RQ2: risks and mitigations}\label{sec:qp2}

This subsection synthesizes the risks and mitigations recorded in the extraction records of the 332 studies of the corpus. The records were extracted by AI and checked by the author in a sample of 67 studies; outside it, they remain without human checking, except for the claims verified one by one against the original. Each study enters once in each category it records. A mitigation is evaluated only when the study compares it with a condition without it, which does not imply effectiveness. The counts of the overview and of Table~\ref{tab:riscos} include all the records; those
of the following subsections come from the partial syntheses, made before the inclusion of S1610, S1617, S667 and of the studies of the last batch of stage 3.

\subsubsection{Overview}

% AFIRMACAO: A45
Functional correctness is recorded in 314 of the 332 studies, followed by cost (156), contamination (138), non-determinism (126), hallucination (99), security (90) and maintainability (41); 275 studies record at least one risk classified as ``other'' (Table~\ref{tab:riscos}).
% AFIRMACAO: A46
Almost all studies (328/332) record at least one mitigation, but only 144/332 evaluate any of them; per task, they are 112/247 in generation, 25/62 in debugging, 12/27 in translation and 7/20 in refactoring.

\begin{table}[!t]
\centering
\footnotesize
\setlength{\tabcolsep}{3pt}
\caption{Distinct studies per risk category and task}
\label{tab:riscos}
\begin{tabular}{p{2.6cm}p{0.75cm}p{0.75cm}p{0.75cm}p{0.75cm}p{0.75cm}}
\hline
Category & Gen. & Deb. & Tr. & Ref. & All \\
\hline
Studies (n) & 247 & 62 & 27 & 20 & 332 \\
\hline
Functional correctness & 232 & 58 & 27 & 18 & 314 \\
Security & 72 & 15 & 6 & 6 & 90 \\
Hallucination & 71 & 20 & 6 & 9 & 99 \\
Contamination & 96 & 31 & 13 & 6 & 138 \\
Non-determinism & 89 & 26 & 13 & 10 & 126 \\
Cost & 119 & 31 & 11 & 6 & 156 \\
Maintainability & 30 & 6 & 3 & 8 & 41 \\
Other & 202 & 54 & 22 & 16 & 275 \\
\hline
With mitigation & 244 & 61 & 27 & 20 & 328 \\
With evaluated mitigation & 112 & 25 & 12 & 7 & 144 \\
\hline
\end{tabular}
\par\smallskip
\parbox{\columnwidth}{\scriptsize Gen.: generation; Deb.: debugging; Tr.: translation; Ref.: refactoring. A study may have more than one task; the task columns do not add up to the total. Count by script from records extracted by AI, checked by sample. ``Evaluated mitigation'': at least one compared with a condition without it.}
\end{table}

\subsubsection{Synthesis by category}

\paragraph{Functional correctness}
% AFIRMACAO: A47
Passing tests is not equivalent to correctness: S1943 reports 62 correct patches among 93 plausible ones (S1943, Table I), S303 reports 186 plausible and 164 correct (S303, Table III) and, in S1562, tests generated by the agent passed wrong code in 90.9\% of the GPT failures (S1562, section 4.5).
% AFIRMACAO: A48
The prompting technique itself can worsen the measured outcome, with records for \textit{few-shot} (S6, section 4.2; S1598, p. 4, pass rate of 0.108 against 0.458 for \textit{zero-shot}), CoT in small models (S319, section 4.3; S1024, Table 5), RAG (S588, Fig. 5a), excess of context or of examples (S275, Figure 5; S1032, Fig. 3) and security preambles (S1693, section V-B).
% AFIRMACAO: A49
Self-correction and repair also make previously correct answers incorrect: S759 reports a $\Delta_{c \to i}$ of 8.4\% in Self-Refine with a small model (S759, Table 1), in S1592, the selection by reward model between the base generation and the hint-guided one corrupted 117 correct solutions and repaired 310 (S1592, section 3.3) and, in S2051, the rate after the first repair iteration stayed below the initial pass@3, measurements that are not equivalent, whereas selecting the best version improves the result (S2051, Table 8).

The mitigations evaluated were execution or compiler feedback loops (S302, S1053, S1077, S1086, S1562), external verification (S1977), interventions on the retrieved context (S670, S887) and control of reasoning according to difficulty (S1133, S1034), part of them by indirect indicators, such as similarity (S1155). Human review and manual checking of plausible patches remain as proposals (S1943, S2057).

\paragraph{Security}
The category brings together vulnerabilities in the generated code, privacy of the data sent to the model and operational security (S2060).
% AFIRMACAO: A50
When measured, vulnerabilities are frequent: in translation, from 28.6\% to 45.0\% of the translations of fixed sources introduce vulnerabilities and from 61.4\% to 71.9\% of the translations of vulnerable sources preserve them (S951, Table IV), and ``improvement'' iterations accumulated 387 distinct vulnerabilities, even with prompts that ask for more security (S739, Tables II and III).

The mitigations evaluated were security instructions in the prompt (S397, S427, S604, S951, S1179, S1195, S1604, S1607, S1693, S1763), injection of security knowledge (S695, S1894), reflection with detector feedback (S1065) and packages without isolation of mechanisms (S391, S2129).
% AFIRMACAO: A51
These evaluations record trade-offs with functional correctness: security guidelines reduced the Pass@1 of Claude-3.7-Sonnet from 16.19\% to 9.52\% (S1195, Table 7), RCI reduced pass@1 on HumanEval (S604, section V-C) and security preambles reduced correctness (S1693, section V-B); in S397, CWE-78 seemed not to be affected by the techniques (S397, section 6.3.2).

Sandbox or container (S137, S1415, S1694), human review between iterations (S739) and local deployment of the model (S1747, S1894) remain as proposals.

\paragraph{Hallucination}
The term has no uniform use; S511, for example, uses it for refactorings that change behaviour.
% AFIRMACAO: A52
A recurring pattern is the reference to non-existent APIs, functions, attributes or identifiers: in S1415, 47 to 48 of 100 reviewed failures are of this nature (S1415, Fig. 5), in S466, 22 cases, in a sample described as of 100 whose categories add up to 110 (S466, section 5.1), and in S109 at least one hardware value was hallucinated in 100\% of the IMU programs with GPT-3.5 (S109, section 5.2).
% AFIRMACAO: A53
Few evaluations measure hallucination itself: in S917, providing the list of signals in the prompt eliminated the non-existent signals observed with the base prompt, in a single task (S917, Tables V to X), whereas in S315 the inclusion of the imports was measured by pass@5 (0.233 without and 0.271 with the instruction, on 10 bugs) and in S1828, S1898 and S1977 the mitigations were measured by Pass@1 or pass rate.

\paragraph{Contamination}
% AFIRMACAO: A54
In the 128 studies that record it in the partial synthesis (138 in the current count, Table~\ref{tab:riscos}), contamination appears in general as a hypothesis; indications appear in S398 (at least 43 occurrences of each HumanEval prompt on GitHub, which indicates exposure, and not inclusion in the training of a specific model; section 3.1), S412 (24 of 98 correct patches lexically identical to the reference, section 6.1.1) and S2057 (122 of 522 bugs in the training data of StarCoder2, PDF p. 27).

The controls adopted (temporal cut-off, new benchmarks, renaming and similarity filters) are not compared with conditions without control.
% AFIRMACAO: A55
Only two studies evaluate a contamination mitigation against a condition without it: S137, by ablation with obfuscated prompts, whose outcome is performance and not contamination (Appendix M), and S2057, which compares trends between Defects4J and Defects4J-clean.

\paragraph{Non-determinism}
Repetitions, seeds and temperature 0 are adopted as a control, in general without comparison.
% AFIRMACAO: A56
Temperature 0 reduces, but does not eliminate, the variation between runs (S124, Table 5; S1695, Table V), and S782 reports a similarity of 0.7 between outputs of the same prompt with temperature 0 (S782, section 4.2); S124 evaluated two prompts: the \textit{concise prompt} made the candidates more deterministic than the base prompt at temperature 0, and CoT increased randomness in that condition (S124, Table 8).

\paragraph{Cost}
% AFIRMACAO: A57
Multi-step techniques add cost when it is measured: MapCoder uses on average 16.7 calls and 21.25 thousand tokens, against one call and 0.64 thousand tokens for the direct prompt (S344, Table 8), the technique of S2140 cost US\$ 0.421 against US\$ 0.017 for Default per 100 problems (S2140, section 7.1) and repair consumed from 5 to 15 times the energy of the single pass (S1415, section 6.1).

The units and the objects measured differ, which prevents aggregating the values; in some studies, the cost is that of the evaluation, and not of the prompting (S892, S1024).
% AFIRMACAO: A58
The cost mitigations evaluated use distinct metrics: candidate pruning reduced the time from 110.1 to 13.2 s per question (S195, Table 10), and routing by difficulty and \textit{adaptive} CoT were measured in tokens and in generation time, respectively (S1133, Table 2; S1034, Table 2).

\paragraph{Maintainability}
% AFIRMACAO: A59
When measured, maintainability gets worse in some scenarios: the C2C condition generated 62 \textit{blocker} or \textit{critical} code smells (56 and 6) against 17 (10 and 7) in the ground truth (S1708, Table VIII), complexity grew with the iterations in association with the vulnerabilities (r = 0.64; S739, section IV-D) and Zero-Verbosity reached a cyclomatic complexity of 54 (S1606, Table I).

Only two evaluated mitigations relate to the topic (S114, S2051).

\paragraph{Topics under ``other''}
In the partial syntheses, among the 248 studies that then recorded risks in the ``other'' category (275 in the current
count), threats to the validity of the studies themselves predominate: validity of the measurement and of the evaluation (100 of these
studies) and external validity (96). There are also context limits, unrequested behaviours (S1008) and human and ethical factors.
% AFIRMACAO: A60
In the latter, S62 reports participants who accepted without scrutiny the reasoning of incorrect patches (S62, section 5.4), and S140 records that \textit{few-shot} raised the social bias (CBS) of GPT-4 from 59.88\% to 68.56\% (S140, Table 7).

\subsubsection{Gaps}

\begin{itemize}
\item \textbf{Risks without evaluation.} Contamination (138 studies in the current count) has two mitigation evaluations in the
partial syntheses, and non-determinism (126), a single prompting one (S124). One possible reading is that the controls work as design care, without verifying the intended effect.
\item \textbf{Outcome different from the risk.} Hallucination mitigations are measured, in general, by Pass@k or pass rate; the term, like ``security'', covers distinct phenomena.
\item \textbf{Trade-offs between risks.} Few studies measure security and functional correctness together, and those that do record losses.
\item \textbf{Packages without isolation.} S1734, S2129 and S2140 do not allow the effect to be attributed to a component of the prompt.
\item \textbf{Uneven coverage.} Translation (27) and refactoring (20) have few studies.
\item \textbf{Checking by sample.} Counts and locators depend on records that, outside the sample of 67 studies, were not checked by a human.
\end{itemize}


\subsection{Quality of the studies}\label{sec:qualidade}
Table~\ref{tab:qualidade} summarizes the quality judgements, made by the extractors. Almost all studies
state a clear objective (327 of 332). Procedure, metrics and reproducibility are partial in almost all of them (312,
324 and 324), which usually reflects the lack of some element, such as the exact model version, the decoding
parameters or the complete prompts. The comparator is adequate in 105 studies and partial in 206. Contamination of the
benchmarks is not addressed in 187 studies, and only 17 address it sufficiently. This is the most
frequent limit to interpreting results on public benchmarks.

\begin{table}[!t]
\caption{Quality of the studies (judgement by AI, $N = 332$)}\label{tab:qualidade}
\centering\footnotesize
\begin{tabular}{@{}lrrrr@{}}
\hline
Criterion & Meets & Partial & Does not meet & Not rep.\\
\hline
Objective & 327 & 5 & 0 & 0\\
Procedure & 20 & 312 & 0 & 0\\
Comparator & 105 & 206 & 21 & 0\\
Metrics & 5 & 324 & 3 & 0\\
Threats to validity & 85 & 216 & 0 & 31\\
Reproducibility & 5 & 324 & 3 & 0\\
Contamination & 17 & 126 & 2 & 187\\
\hline
\end{tabular}
\end{table}

%% QP3
% English translation of capitulos/04c-qp3.tex, made by Claude Code (Claude Fable 5.2) on 24 November 2026.
% Pending review by the author. When the Portuguese file changes, update this one (scripts/verificar_traducao.py).
\subsection{RQ3: taxonomy and decision-support table}\label{sec:qp3}

This subsection brings together the practical contribution of the work: a taxonomy of the techniques
extracted and a decision-support table, in the version revised after the examination in the scenarios. Both are based on 332 records
and 805 comparisons extracted by AI and checked by the author in a sample (Section~\ref{sec:conferencia}). The counts come from
\texttt{scripts/sintetizar.py}, and the rows of the table, from the partial syntheses.

\subsubsection{Taxonomy}

The taxonomy is multi-label: a study may have more than one task, and a technique, more
than one mechanism. The records register five dimensions:
\begin{itemize}
  \item \textit{Task}: generation (247 studies), debugging (62), translation (27) and
    refactoring (20); 22 studies have more than one task.
  \item \textit{Mechanism}: eleven labels (instruction, context, constraints, examples,
    decomposition, reasoning, feedback, sampling, multi-agent, prompt optimization and
    other). The record separates the complete package, with all the mechanisms present in the prompt
    of the technique, from the mechanisms that the technique adds to the base prompt;
    Table~\ref{tab:mecanismos} counts only those added.
  \item \textit{Role}: focal, when the technique is the object of the study, or comparator, when
    it serves as a base or as an alternative to another technique.
  \item \textit{Type of counterfactual}: package (whole technique against the condition without
    it), ingredient (an isolated component), paired contrast, cumulative (successive
    rounds, without a parallel condition) or none. Among the focal techniques, there are 373 of
    package, 212 of ingredient, 3 of paired contrast, 13 cumulative and 111 without a
    counterfactual (unit: technique). Counting also the components of the techniques, 191 of the
    332 studies have at least one isolated ingredient or paired contrast.
  \item \textit{Conditions}: factors on which the results depend, according to the syntheses.
\end{itemize}

The syntheses point to the following recurring conditions:
\begin{itemize}
  \item available executor, tests or compiler, on which feedback and
    candidate selection depend;
  \item number of calls and of candidates, rarely equalized between the conditions;
  \item size and family of the model, with sign reversals within the same study;
  \item level of the task: function, class or repository;
  \item oracle information, such as perfect localization of the defect or APIs and solutions
    taken from the reference;
  \item components outside the scope of prompting (fine-tuning, trained retrievers,
    model selection), recorded in 303 of 332 studies;
  \item metric (textual similarity, tests or scanners) and domain, such as hardware and
    domain languages, in general with their own sets.
\end{itemize}

\subsubsection{Decision-support table}

Version 1 of the table consolidated the rows of the partial syntheses into 22, one per grouping
and task, and was frozen for the examination in the scenarios (Section~\ref{sec:exame}). After the
examination, the author decided to revise it. Version 2 (Tables~\ref{tab:apoio} and~\ref{tab:apoiob})
uses the same syntheses and has 31 rows, with five changes:
\begin{itemize}
  \item feedback was separated by type of verification, tests or compiler and
  analyser, in generation (G1 and G2) and in translation (T1 and T2);
  \item each row states the measured result, such as compilation, passing the tests,
  correctness of the repair or complexity;
  \item the condition makes explicit when the evidence assumes perfect localization or a test that
  reproduces the defect (D1, D3 and D6);
  \item the refactoring row was divided into specified type, API knowledge and
  examples (R1 to R3);
  \item each indication is marked as favourable, contrary, mixed, indirect or
  insufficient. Five new rows declare insufficient evidence: the distribution of a
  fixed budget of calls (G13), domain languages and small models (G14), repair
  without reproduction (D7), legacy language (T5) and repository scale (R5). The separations
  created G2, with indirect evidence, T2, insufficient, and R1 to R3.
\end{itemize}
In the mixed, contrary or insufficient rows, the guidance is not to recommend the technique
as a default. The strength is descriptive, without a score, and repeats the count of v1 when a row was
divided. v1 and v2 are in the repository (\path{evidencias/sinteses/}). The findings
below support the central rows.

In execution feedback in generation, the budget matters.
% AFIRMACAO: A91
With the number of attempts equalized, the multi-turn variants with feedback stayed at
2.4\% and 2.8\% of pass 1@3, against 2.9\% without iteration (Llama 3.0 8B), and at 6.9\% against
7.7\% (Llama 3.1 8B); SEIDR solved on average 11.5 problems of PSB2, against 12.2 for
sampling without repair (S479-C2; S479-C3; S479-C4; S645-C1).
% AFIRMACAO: A92
Without this equalization, five comparisons that isolate the loop and have absolute
values recorded a more favourable outcome with it, in the sense of each metric, and removing the iterative verification from CG-CG
reduced Pass@1 from 80.3\% to 70.1\% (S878-C1; S1047-C3; S1053-C3; S1077-C1;
S1086-C6; S1977-C2).
% AFIRMACAO: A93
In multi-agent systems, the gain comes with more calls: MapCoder and
LLMDebugger obtained 64.02\% and 67.07\% pass@1, against 39.63\% for \textit{zero-shot}, with 8.1
and 3.4 calls on average, against 1.0 (S1747-C3; S1747-C4).

% AFIRMACAO: A94
The reasoning instruction added to the generation prompt had a mixed result: in
12 comparisons from 12 studies, under heterogeneous conditions and metrics, 5 record an improvement and 7, a deterioration, a
descriptive count that is not equivalent to 12 isolated effects of the instruction (S124-C2; S207-C1; S309-C2; S319-C1; S491-C3; S508-C1;
S560-C2; S569-C1; S1024-C4; S1121-C1; S1229-C2; S1606-C4).
% AFIRMACAO: A95
Self-correction without execution, with Qwen2.5-1.5B, stayed below Simple-prompting: 38.2\%
(Self-Refine) and 38.8\% (SETS), against 44.6\% Accuracy@t2 (S759-C1; S759-C3).
% AFIRMACAO: A96
In the security instructions, 13 variants without iteration raised the vulnerable samples
of GPT-3.5-turbo from 5.74\% to values between 5.99\% and 7.82\% (S604-C1 to S604-C10;
S604-C13 to S604-C15), whereas, with GPT-3, the instructional variants reduced the
weaknesses counted by Bandit from 78 to values between 47 and 72 (S397-C2 to S397-C6).
Since model, metric and set differ, the two studies do not form a direct contrast.
% AFIRMACAO: A97
Oracle information produced the largest gains among the contrasts cited, without representing real use: the file
context added to the APIs taken from the reference raised the Pass@1 of GPT-4o Mini on RepoExec from 16.62\% to 37.75\%
(S670-C5), and the official solution raised the pass@1 of DeepSeek-Coder-V2 on MuPA from 11.84 to 38.79, the highest value among the conditions cited, whereas the self-generated hints stayed between 11.92 and 13.21 (S2145-C1; S2145-C3;
S2145-C5; S2145-C6).

% AFIRMACAO: A98
In repair, explicit reasoning worsened the result in 5 of 6 comparisons from 5
studies; the exception, the \textit{few-shot} CoT of ThinkRepair, also changes the selection of examples
(S275-C2; S354-C2; S1429-C1; S1429-C2; S1598-C2; S412-C2).
% AFIRMACAO: A99
Structural or bulky diagnostic information also made things worse: with Qwen2.5-Coder:7B on
QuixBugs, PDG and test file reduced the repair rate from 62.5\% to 55\% and
57.5\%, and, on Defects4J v1.2, prompts with stack trace localized 32 and 41 bugs, against 47
without it (S1674-C3; S1674-C4; S1678-C1; S1678-C2).
% AFIRMACAO: A100
Three repair comparisons with fixed examples recorded no improvement: the
overall\_score fell from 0.783 to 0.643, the plausible patches fell from 205 to 200, and
the failures fixed stayed at 11 against 11 (S1598-C1; S2057-C1; S1565-C4).

% AFIRMACAO: A101
In translation, explaining before translating reduced pass@1 from Python to JavaScript with
GPT-3.5 from 85.5\% to 84.8\%, 84.3\% and 83.2\%, and the two-step approaches reduced
the success rate on CodeNet from 28.6\% to values between 2.6\% and 11.2\% (S193-C1;
S193-C2; S193-C3; S836-C2; S836-C3; S836-C5; S836-C6).
% AFIRMACAO: A102
In refactoring, specifying the type of refactoring raised the recall of GPT-4 from 15.6\% to
52.2\%, and the TerraMod package raised the Pass@1 of Mistral-Large from 4\% to 76\%, both
from low baselines (S511-C2; S2122-C1).

% English translation of capitulos/04c2-tabela-apoio-v2.tex, made by Claude Code (Claude Fable 5.2) on 24 November 2026.
% The Portuguese file is generated by scripts/gerar_tabela_apoio.py from evidencias/sinteses/tabela-apoio-v2.csv;
% this file translates the generated .tex and is not regenerated by the script. Pending review by the author.
\providecommand*{\rr}{\raggedright\hspace{0pt}}
\begin{table*}[!tp]
\caption{Decision-support table, version 2: generation (G1--G12) and common rows (G13--G14)}\label{tab:apoio}
\centering
\fontsize{6.6}{7.4}\selectfont
\renewcommand{\arraystretch}{0.95}
\setlength{\tabcolsep}{2.5pt}
\renewcommand*{\rr}{\raggedright\hspace{0pt}}
\begin{tabular}{p{0.55cm}p{2.0cm}p{2.4cm}p{1.9cm}p{5.6cm}p{1.6cm}p{3.1cm}}
\hline
\rr\textbf{ID} & \rr\textbf{Grouping} & \rr\textbf{Condition and verification} & \rr\textbf{Measured result} & \rr\textbf{Evidence and indication} & \rr\textbf{Strength} & \rr\textbf{Caveats}\tabularnewline
\hline
\rr G1 & \rr Test or execution feedback in a loop & \rr Executable tests or executor; single-call base & \rr Pass@1; passing the tests & \rr \textit{Favourable; contrary with equalized budget.} Higher value with the loop, in a package (S150-C1) and isolated (S150-C2; S1977-C2). In the conditions extracted with the number of attempts equalized, no gain over sampling (S479-C2; S645-C1) & \rr approx. 37 (pkg. approx. 21; ingr. approx. 15), count of v1 for all types of feedback & \rr Calls not equalized in most; other ingredients in the package; passing the tests does not guarantee the requirement\tabularnewline
\rr G2 & \rr Compiler or static analyser feedback & \rr Compiler, IDE or security analyser, without functional tests & \rr Compilation; secure generation accuracy; validity & \rr \textit{Indirect.} Few studies and weak designs: compilable from 8 to 17 in cumulative rounds, without a parallel condition (Solidity; S2055-C1); security feedback in a package with a 32B model (S1065-C1); static report tied with self-review (S1763-C3) & \rr 3 studies & \rr Compiling does not measure functional correctness; a scanner does not demonstrate security\tabularnewline
\rr G3 & \rr Self-critique without execution & \rr The model itself critiques and rewrites, without executing & \rr Weaknesses flagged by scanner; accuracy; completion rate & \rr \textit{Mixed.} Fewer weaknesses with RCI (S397-C1; S604-C11); deterioration with a 1.5B model in one study (S759-C1); small gain (S1204-C2). Do not recommend as a default & \rr 6 (ingr. 1) & \rr One deterioration with 1.5B does not establish the behaviour of all small models; statistical null in S397\tabularnewline
\rr G4 & \rr Multi-agent with roles & \rr Tester, reviewer or evaluator; more calls, in a number that varies per arrangement & \rr Pass@1; compilation (S492) & \rr \textit{Favourable to the package.} Gains against a direct base (S344-C1; S2131-C1) and in ablations (S448-C3); deterioration in compilation in S492-C1 & \rr approx. 29 (pkg. approx. 29; ablation approx. 11) & \rr Arrangement against single agent only in S530; MapCoder with 8.1 calls on average (S1747-C3); fine-tuning of agents\tabularnewline
\rr G5 & \rr Sampling and selection & \rr Several candidates; selection by generated tests or value model & \rr Pass@1; creativity (S399) & \rr \textit{Favourable to selection; contrary to simple sampling.} Improvement in selection ablations (S416-C2; S724-C2; S1034-C2); deterioration with simple sampling (S399-C3) & \rr approx. 17 (ingr. approx. 9) & \rr Selection by the evaluation oracle is an upper bound (S399); trained components (S1034, S1592)\tabularnewline
\rr G6 & \rr Reasoning instruction (CoT) & \rr The same model generates the reasoning & \rr Pass@1 & \rr \textit{Mixed.} 5 improvements and 7 deteriorations in 12 comparisons (S319-C1; S1229-C2; S569-C1; S1121-C1). Do not recommend as a default & \rr 12 (ingr. 11; pkg. 1) & \rr Non-significant null (S560); token ceiling (S1606); reversals between models\tabularnewline
\rr G7 & \rr Planning or decomposition before the code & \rr Two or more steps; function level & \rr Pass@1 & \rr \textit{Favourable to the package.} Improvement in a package (S1-C1) and in ablations (S1133-C3); deteriorations with a step generated by a small model (S216-C2) and with knowledge generation (S2075-C2) & \rr approx. 20 (pkg. approx. 16; ingr. approx. 10) & \rr More calls; cost not measured; in a Java class, holistic generation outperformed the decomposed one (S365)\tabularnewline
\rr G8 & \rr Retrieved context (API, file, repository) & \rr Accessible repository or documentation & \rr Pass@1; accuracy; BLEU (S1951) & \rr \textit{Favourable; contrary in part.} Improvement in a package and in ablations (S357-C1; S852-C2); nulls and reversals (S806-C1; S1951-C1) & \rr approx. 45 (ingr. approx. 20) & \rr Oracle context is an upper bound (S670, S2145); more tokens; trained retrievers\tabularnewline
\rr G9 & \rr In-context examples & \rr Fixed, random or retrieved demonstrations & \rr Pass@1 & \rr \textit{Mixed.} Small differences of variable sign when examples are merely added (S432-C1; S591-C1); larger with trained selection (S170-C1); deterioration with faulty samples (S2075-C1) & \rr approx. 24 (ingr. approx. 14) & \rr Reversals by model and temperature; components outside the scope\tabularnewline
\rr G10 & \rr Specification and tests in the prompt & \rr Tests, pre- and postconditions or format rules in the prompt & \rr Accuracy; Pass@1; assertions passed & \rr \textit{Favourable in part.} Improvement with tests in the prompt (S1747-C1; S1086-C5) and with pre- and postconditions (S1812-C1); deterioration with Zero-Verbosity (S1606-C3) & \rr approx. 25 (ingr. approx. 9) & \rr Tests exposed in the prompt may also evaluate (S1613); S1812 not significant and aggregated\tabularnewline
\rr G11 & \rr Security instructions and constraints & \rr Prefixes, persona, CWE or guidelines & \rr Vulnerable samples by scanner; exposure to risk patterns; Pass@1 & \rr \textit{Mixed.} Deterioration and improvement, on distinct metrics (S604-C1; S397-C2; S1607-C1); lower Pass@1 with guidelines (S1195-C4). Do not recommend as an isolated measure & \rr 7 (pkg. and ingr.) & \rr Benchmark and scanner do not demonstrate security\tabularnewline
\rr G12 & \rr Wording, language, compression and automatic optimization & \rr Wording, natural language, pruning or search of the prompt & \rr Pass@1; incorrect solutions & \rr \textit{Mixed; insufficient per variant.} Rewordings worsen or do not change (S2162-C1; S323-C2); a language other than English worsens (S495-C1); guided pruning improves (S1644-C1). Do not recommend without a local test & \rr approx. 20 (ingr. approx. 8) & \rr One model or one study per variant; optimization in few studies (S650; S2053)\tabularnewline
\rr G13 & \rr Distribution of a fixed budget of calls (all) & \rr Choosing between repairing, generating new candidates or planning, with a limit of calls & \rr Pass 1@3; problems solved & \rr \textit{Insufficient.} The corpus does not guide how to split the calls. In the conditions extracted from the two studies with the budget of attempts equalized, the feedback loop did not outperform sampling (S479-C2; S645-C1); S479 reports gains with other models and budgets & \rr 2 studies & \rr 8B models and GPT-3.5 on competitive problems; they equalize attempts, not tokens or cost\tabularnewline
\rr G14 & \rr Domain languages and small local models & \rr E.g.: Solidity; models up to 8B & \rr Compilation; CodeBLEU (S556); Pass@1 & \rr \textit{Insufficient.} No contrast of a prompting technique in Solidity (S2055-C1 is cumulative; S556 only measures CodeBLEU). With small models, rows G3 and G6 record deteriorations (S759-C1; S1121-C1) & \rr no count of its own & \rr The evidence of the other rows comes from other models and languages\tabularnewline
\hline
\end{tabular}
\par\smallskip
\parbox{\textwidth}{\scriptsize Evidence: favourable, contrary or mixed, according to the direction of the comparisons; indirect, when it comes from another condition, domain or metric; insufficient, when the corpus has no contrast for the condition. Strength: comparisons and studies, without a score. Extraction by AI, checked by sample; syntheses of the 302 records used in v1.}
\end{table*}

\begin{table*}[!tp]
\caption{Decision-support table, version 2: debugging (D), translation (T) and refactoring (R)}\label{tab:apoiob}
\centering
\fontsize{6.6}{7.4}\selectfont
\renewcommand{\arraystretch}{0.95}
\setlength{\tabcolsep}{2.5pt}
\renewcommand*{\rr}{\raggedright\hspace{0pt}}
\begin{tabular}{p{0.55cm}p{2.0cm}p{2.4cm}p{1.9cm}p{5.6cm}p{1.6cm}p{3.1cm}}
\hline
\rr\textbf{ID} & \rr\textbf{Grouping} & \rr\textbf{Condition and verification} & \rr\textbf{Measured result} & \rr\textbf{Evidence and indication} & \rr\textbf{Strength} & \rr\textbf{Caveats}\tabularnewline
\hline
\rr D1 & \rr Test or execution feedback in a loop & \rr Tests that reproduce the defect, debugger or counterexample & \rr Correct fixes; plausible patches (S62) & \rr \textit{Favourable.} Improvement in most (S62-C2; S412-C3); tie with a 1B model (S114-C2) & \rr 13 (pkg. 8; ingr. 4) & \rr Assumes a test that triggers the defect; calls not equalized; hints from the answer key (S65); drop on Codeforces (S2164); a plausible patch is not a verified fix\tabularnewline
\rr D2 & \rr Diagnostic information in the prompt & \rr Failing test with error message, test output or defect report & \rr Correct fixes; repair rate; bugs localized (S1678) & \rr \textit{Favourable; contrary in part.} Improvement with triggering test and message (S647-C2), defect report (S647-C3) and test output (S1674-C5); deterioration with PDG (S1674-C3) & \rr 20 (ingr. in at least 8) & \rr Measured with perfect localization in S647; the deterioration with stack trace (S1678-C1) is in the localization of the defect, not in the repair; varies with the type of bug\tabularnewline
\rr D3 & \rr Retrieved context and ingredients & \rr Repository, history or base of patterns; in general with perfect localization or triggering tests & \rr Correct fixes; repair rate; exact match & \rr \textit{Favourable to the package.} Improvement in a package and in ablations (S275-C1; S435-C3; S2046-C2); injecting all the context at once worsens (S821-C2) & \rr 15 (ingr. 7) & \rr Perfect localization (S2046); triggering tests (S303); unequal budgets; different models between conditions\tabularnewline
\rr D4 & \rr Examples & \rr Fixed, historical or selected & \rr Exact match; programs fixed; overall score & \rr \textit{Mixed.} Improvement with selection or two examples (S25-C2; S53-C1); deterioration or tie with fixed or reference examples (S552-C3; S1598-C1). Do not recommend fixed examples & \rr 8 (ingr. in at least 4) & \rr Reversal between models (S2057); diverging tables (S25)\tabularnewline
\rr D5 & \rr Explicit reasoning & \rr CoT, ReAct or explanation of the error as the only mechanism added & \rr Repair rate; Pass@1; overall score & \rr \textit{Contrary.} Deterioration in 5 of 6 comparisons in which reasoning is the only mechanism added (S275-C2; S354-C2; S1429-C1; S1598-C2). Do not adopt as a default on the basis of this corpus; evaluate under the conditions of use & \rr 5 (ingr. 3; pkg. 2) & \rr The improvement (S412-C2) includes selection of examples; few contrasts, with a 1.3B model (S1598) and the quantum domain (S1429)\tabularnewline
\rr D6 & \rr Decomposition, planning and agents & \rr Several stages, tools or debate; in general with execution that reproduces the defect & \rr Correct fixes; repair rate & \rr \textit{Favourable to the package.} Improvement in a package (S810-C1; S699-C1) and in ablations (S810-C2; S1304-C3) & \rr 12 (pkg. 12; ablation in 4) & \rr Execution traces require reproduction (S1304-C2); different calls and models; reversals in subtotals (S303, S1304)\tabularnewline
\rr D7 & \rr Repair without reproduction and with unknown localization & \rr Only the defect report; the suite does not trigger the failure & \rr Correct fixes; plausible patches & \rr \textit{Insufficient.} There is no evidence to recommend the complete process. The defect report in the prompt is the only applicable favourable element, measured with perfect localization (S647-C3). Compilation and existing tests verify regressions, not the reported defect & \rr no count of its own & \rr Defects4J provides triggering tests; S303 and S2046 assume tests or localization\tabularnewline
\rr T1 & \rr Repair with tests and compiler & \rr Executable tests and compiler; iterative repair & \rr Computational accuracy; pass rate & \rr \textit{Favourable; contrary in one study.} Improvement in most: package with repair (S732-C1); tests in the prompt and repair (S329-C2; S329-C3). Drop in S2051-C1 (non-equivalent metrics); null at the ceiling (S919-C3) & \rr 11 (pkg. 9; ingr. 2) & \rr In S329, the tests in the prompt accompany most of the difference; S732 is a package with its own repair steps and limits; cost not measured\tabularnewline
\rr T2 & \rr Compiler-only feedback & \rr Compiler, without tests & \rr Compilation & \rr \textit{Insufficient.} Compiling does not demonstrate equivalence of behaviour; in S1163-C1, compilation stayed at the ceiling in both conditions & \rr 1 study & \rr No comparison measures equivalence with only the compiler\tabularnewline
\rr T3 & \rr Examples & \rr k examples; structural or retrieved selection & \rr Computational accuracy; BLEU; compilation & \rr \textit{Mixed.} Small improvement (S674-C1; S1032-C1) or deterioration (S964-C1). Do not recommend as a default & \rr 6 (ingr. 3; pkg. 3) & \rr BLEU does not measure correctness; the effect does not grow with k (S409)\tabularnewline
\rr T4 & \rr Explanation and two steps & \rr Explaining or representing before translating & \rr pass@1; rate of successful translation & \rr \textit{Contrary.} Deterioration (S193-C1; S836-C2; S836-C3); one-shot with reasoning improves (S836-C1). Do not adopt as a default on the basis of this corpus; evaluate under the conditions of use & \rr 2 (ingr. 1; pkg. 1) & \rr S836-C1 confounds example and reasoning; one model per study\tabularnewline
\rr T5 & \rr Legacy language or platform migration & \rr E.g.: COBOL; change of version or of platform & \rr — & \rr \textit{Insufficient.} No study of the corpus translates COBOL; migrations only appear without a counterfactual or in a cumulative contrast (S1220; S1803; S1836) & \rr no count of its own & \rr The absence of evidence does not demonstrate that the techniques are ineffective\tabularnewline
\rr R1 & \rr Refactoring type specified & \rr Expected type indicated in the prompt & \rr Recall in the identification of opportunities & \rr \textit{Indirect.} Recall from 15.6\% to 52.2\% (S511-C2) & \rr 1 study & \rr Measures the identification of opportunities, not the refactoring performed\tabularnewline
\rr R2 & \rr API knowledge: schema and changelog & \rr API documentation and changelog in the prompt & \rr Pass@1 & \rr \textit{Favourable in one study; indirect outside it.} Package from 4\% to 76\%; without the schema, 56\%; without the changelogs, 60\% (S2122-C1 to S2122-C3) & \rr 1 study & \rr Single study, on Terraform; number of calls not reported\tabularnewline
\rr R3 & \rr Examples & \rr Selected or one-shot examples; low baseline & \rr Passing the tests; Pass@1; correctness (S1822); complexity (S202) & \rr \textit{Favourable in pass rate; mixed in complexity.} Improvement (S202-C1; S4-C2; S1822-C1); S202 also measures complexity, with small differences that depend on the condition & \rr count of v1 for the three refactoring ingredients together: 11 (ingr. 6) & \rr Own sets; passing the tests does not imply reducing complexity; number of examples not established\tabularnewline
\rr R4 & \rr Feedback, iteration and reasoning & \rr Tests, compilation or detector; CoT & \rr Proportion that passes the tests; compilation; syntactic validity & \rr \textit{Insufficient.} Improvements (S1054-C1; S2183-C1); higher compilation without the feedback (S1685-C3); CoT with a small gain (S4-C1). Do not recommend as a default & \rr 7 (pkg. or ingr. in 6) & \rr Comparators with several ingredients; S739 without a baseline\tabularnewline
\rr R5 & \rr Repository scale and performance & \rr Interdependent changes in many files; performance requirement & \rr — & \rr \textit{Insufficient.} No study with a counterfactual at this scale; performance only appears against human code (S958) & \rr no count of its own & \rr Per-file checks do not establish the success of the complete migration\tabularnewline
\hline
\end{tabular}
\par\smallskip
\parbox{\textwidth}{\scriptsize Evidence: favourable, contrary or mixed, according to the direction of the comparisons; indirect, when it comes from another condition, domain or metric; insufficient, when the corpus has no contrast for the condition. Strength: comparisons and studies, without a score. Extraction by AI, checked by sample; syntheses of the 302 records used in v1.}
\end{table*}


\subsubsection{How to read the table and what it does not do}

The table is read by row: given the task and the condition of use, it summarizes the direction of the
extracted comparisons, the strength with which they were obtained and the caveats that prevent
generalization. ``Favourable'' indicates that most comparisons of the grouping recorded an
improvement on the metric of each study, and not that the technique will work in another project.
When the row is favourable to the package, the gain cannot be attributed to a single mechanism. The
table does not rank techniques: benchmarks, metrics and models differ between the rows, and
no average was computed. It also does not measure effectiveness in real use, since benchmarks do not
demonstrate productivity, security or cost reduction. The proportion of improvements among
the 805 comparisons (516 improvements, 145 deteriorations, 28 equal, 91 without direction and 25
descriptive) reflects the choice of the main comparison by the authors and publication
bias, and does not estimate an effect. Groupings with a single study, such as prompt
optimization and security requirements in translation (S650, S951), were left out of the table and
count as insufficient evidence. The two versions use the syntheses of 302 records; S1610, S1617, S667 and the studies of the last
batch of stage 3, extracted later, do not enter its rows. S433, cited in a row of v1, was later grouped
with S2053 for being a version of the same work (X1), and v2 cites S2053. Finally, the table reflects extraction by AI checked only by
sample: counts and directions may change with corrections outside it.

% English translation of capitulos/04e-exame.tex, made by Claude Code (Claude Fable 5.2) on 24 November 2026.
% Pending review by the author. When the Portuguese file changes, update this one (scripts/verificar_traducao.py).
\subsection{Examination of the table in the scenarios}\label{sec:exame}
Table~\ref{tab:cenarios} summarizes the eight synthetic scenarios, two per task. Each one fixes language and domain,
model, verification resources and call budget. The evidence expectations are hypotheses to be confronted. The
scenarios and decision-support table v1 were frozen with SHA-256 before the examination. For each scenario, the examination applies the
rows of the table, checks the studies cited in the records and registers limitations and ambiguities. Claude Code drafted the examination and a proposed judgement, and the author made the final judgement. The examination does not measure effectiveness.

\begin{table}[!t]
\caption{Scenarios used in the examination of version 1 of the decision-support table}\label{tab:cenarios}
\centering\scriptsize
\setlength{\tabcolsep}{3pt}
\begin{tabular}{@{}p{0.55cm}p{3.2cm}p{2.4cm}p{1.3cm}@{}}
\hline
ID & Situation & Verification & Budget\\
\hline
G1 & Python functions from a docstring; large model & Public unit tests & 5 calls\\
G2 & Solidity 0.8 contracts; local model up to 8B & \texttt{solc} and static analyser & 3 calls\\
D1 & Defect in a Python function & Failing test and suite & 10 calls\\
D2 & Defect in a Java repository from a report & Compilation; suite does not reproduce & 20 calls\\
T1 & C99 module to safe Rust & \texttt{rustc} and ported tests & 5 rounds\\
T2 & Batch COBOL to Java 17 & Java compiler only & 3 calls\\
R1 & Long Python methods, lower complexity & Tests and complexity meter & 3 calls\\
R2 & API migration in about 200 Java files & Compilation, tests, benchmarks, changelog & 2 per file\\
\hline
\end{tabular}
\end{table}

The author judged the eight scenarios on 17 November 2026 (Table~\ref{tab:julgamento}). On the three criteria, the result was:
\begin{itemize}
  \item clarity: 1 meets and 7 partial;
  \item grounding in the sources: 2 meet, 5 partial and 1 does not meet;
  \item concrete guidance: 1 meets, 6 partial and 1 does not meet.
\end{itemize}
The author agreed with 22 of the 24 judgements proposed by Claude Code and downgraded to partial the grounding of
T1 and R1. In T1, the result cited (S732-C1) comes from a package with several repair steps and limits per step, which
is not equivalent to five total rounds per function. In R1, the highlighted results do not support at the same time the
reduction of complexity and that of smells.

\begin{table}[!t]
\caption{Author's judgement per scenario and criterion}\label{tab:julgamento}
\centering\scriptsize
\begin{tabular}{@{}lccc@{}}
\hline
Scenario & Clarity & Grounding & Guidance\\
\hline
G1 & partial & meets & partial\\
G2 & partial & partial & partial\\
D1 & partial & meets & meets\\
D2 & partial & partial & does not meet\\
T1 & meets & partial & partial\\
T2 & partial & does not meet & partial\\
R1 & partial & partial & partial\\
R2 & partial & partial & partial\\
\hline
\end{tabular}
\end{table}

D1 is the only scenario in which grounding and guidance meet the criteria: the table leads to an executable script, with the
function, the test and the error message in the prompt and verification by the suite. The guidance fails in D2, in which the
suite does not reproduce the defect and the localization is unknown, and the grounding fails in T2, in which there is only a compiler
for a COBOL migration that requires equivalence of behaviour. In both cases, the table does not declare that the
evidence is insufficient. The ambiguities that the author recorded recur across the scenarios:
\begin{itemize}
  \item the table does not say how to split a fixed budget of calls between repair and new candidates (G1, G2 and T1);
  \item the feedback condition puts together compiler, tests and static analyser and does not distinguish the result
  measured, such as compilation, passing the tests or fixing the defect (G2, T1 and T2);
  \item passing the suite does not guarantee that the change fixes the intended defect or that it preserves the behaviour
  (D1, D2 and T2);
  \item the dependence on known localization or on a test that reproduces the defect only appears in the caveats (D2);
  \item the refactoring row puts together specification, domain knowledge and examples, which have different foundations
  (R1 and R2).
\end{itemize}
The author also corrected the derived guidance in some scenarios. In G2, the appropriate wording is not to recommend
self-critique without execution as a default, since one deterioration with a 1.5B model does not establish the behaviour of all
models up to 8B. In D2, compilation and existing tests still give feedback on the change, although they do not verify
the reported defect. In T1, the study cited does not have a budget equal to that of the scenario. In R1, from one to three examples is not
a range established by the evidence.

The author decided to apply the revisions in a version 2 of the table:
\begin{itemize}
  \item distinguish the types of feedback and make explicit the result measured (compilation, passing the tests,
  correctness of the repair, complexity or performance);
  \item make explicit the dependences on localization and on reproduction of the defect;
  \item declare the conditions without sufficient evidence and the cases in which the corpus does not guide the distribution of the
  calls;
  \item separate the ingredients grouped in the refactoring row;
  \item differentiate absence of evidence, contrary result and indirect evidence.
\end{itemize}
In the first proposal, the eight hypotheses of the scenarios were given as confirmed, H-G1 only in part. The author asked
that each hypothesis be judged separately, admitting partial confirmation or an inconclusive result. The examination assesses the
guidance offered by the table, and not the effectiveness of the techniques. The complete examination is in
\path{evidencias/sinteses/exame-cenarios-v1.md}, and the author's judgements, in
\path{evidencias/sinteses/exame/julgamentos-autor.json}.
Version 2 of the table, revised after the examination, is in Section~\ref{sec:qp3}. In the separate judgement, the author confirmed six hypotheses and confirmed two in part
(\path{evidencias/sinteses/exame/hipoteses-proposta.md}). H-G1 is confirmed in part because there is evidence to
consider tests, feedback and planning, but not to choose between them with five calls. H-T1 is
confirmed in part because the result cited comes from a package with its own steps and limits and does not cover the requirement
of safe Rust. No hypothesis was inconclusive. The six confirmed include the four that predicted insufficient
evidence (H-G2, H-D2, H-T2 and H-R2): in these scenarios, not recommending a technique as a default is the appropriate
answer.


\subsection{Status of the human check}\label{sec:conferencia}
% English translation of capitulos/04d-conferencia.tex, made by Claude Code (Claude Fable 5.2) on 24 November 2026.
% Pending review by the author. When the Portuguese file changes, update this one (scripts/verificar_traducao.py).
The check sample has 67 studies (Section~\ref{sec:extracao}). The author checked all of them on 17 November 2026, with
about 6 hours of work in total for the first 66, according to the author's account. S1935 entered the sample later, in
place of S1721, excluded from the corpus, and was checked next. In each study, the author judged as correct or to be
corrected each block of the record (classification, context, risks, heterogeneity, quotations, quality and each
comparison), with the right value and the page.

The record has two levels. In 11 studies, the judgement of each block is stored in the form. In the other 56, the
author declared that all the blocks were correct, without storing them one by one. The proportion below treats these 56
studies as without correction, according to the declaration:
\begin{itemize}
  \item studies with some correction: 6 of 67 (5 of the 61 extracted by Opus and 1 of the 6 extracted by Codex);
  \item blocks corrected: 18 of 571 (3.2\%), 17 of 521 in the Opus records and 1 of 50 in the Codex ones.
\end{itemize}

The corrections were of five types:
\begin{itemize}
  \item values that the extractor judged illegible in a low-resolution table and that the author read with magnification
  (S1743 and S650);
  \item the selection of the comparisons, with a variant proposed by the authors of the study treated as a focal technique (S77);
  \item the classification of mechanisms and the unit of a metric not declared by the study (S1935);
  \item outdated justifications or incorrect generalizations in the heterogeneity (S6 and S1935);
  \item a word of a quotation transcribed differently from the printed text (S2090).
\end{itemize}
No correction reversed the direction of a comparison already extracted. Those of S1743 added four results that
were missing, and that of S77 replaced an ablation with a focal variant. The changes are recorded field by
field, and the raw outputs of the extractors were not modified. With 6 Codex studies in the sample, it is not possible
to compare the extractors.

The author also checked against the original the 120 empirical claims in the text, each with the passage in which it appears and the
comparisons or the reference that support it. The author judged 89 as supported, 30 as to be corrected and 1 as not
verifiable. The corrections were of framing and wording, and not of direction: indicating the subtable, the category or the
condition the number comes from (for example, the split of CoderEval in S302 or the subset of LeetCode in S2164),
distinguishing a non-significant result from absence of effect, separating cumulative stages (S511) and non-equivalent
measurements (S2051), and removing a debugging comparison from a generation count (S654-C2). The sentences were rewritten according to these indications,
including the third denominator of S1880 (206 valid codes with RCI and RAG), which was missing in A106. The non-verifiable claim (A84) depended on a fraction calculated in the
partial syntheses, which the author could not reproduce; the sentence now cites only the checked example.
The claims remain with draft status in the data, with the reviewer recorded in the supported ones, because the
``supported'' status also requires checked metadata, which are left for the final version.

